\documentclass[11pt]{amsart}

\usepackage[utf8]{inputenc}
\usepackage[T1]{fontenc}
\usepackage{lmodern}
\usepackage{amsmath,amssymb,amsthm,mathtools}
\usepackage[margin=1.1in]{geometry}
\usepackage{enumitem}
\usepackage{booktabs}
\usepackage{array}
\usepackage{xcolor}
\usepackage[colorlinks=true,linkcolor=blue!50!black,citecolor=green!40!black,urlcolor=blue!60!black]{hyperref}
\usepackage[capitalize,nameinlink]{cleveref}
\emergencystretch=3em

\newtheorem{theorem}{Theorem}[section]
\newtheorem{lemma}[theorem]{Lemma}
\newtheorem{proposition}[theorem]{Proposition}
\newtheorem{corollary}[theorem]{Corollary}
\newtheorem{conjecture}[theorem]{Conjecture}
\newtheorem{question}[theorem]{Question}
\newtheorem*{mainthm}{Theorem 1}
\theoremstyle{definition}
\newtheorem{definition}[theorem]{Definition}
\theoremstyle{remark}
\newtheorem{remark}[theorem]{Remark}

\newcommand{\R}{\mathbb{R}}
\newcommand{\Dt}{\Delta_2}
\newcommand{\mult}{\mu}
\newcommand{\eps}{\varepsilon}
\newcommand{\conv}{\operatorname{conv}}
\newcommand{\turn}{\operatorname{turn}}
\newcommand{\dist}{\operatorname{dist}}
\newcommand{\Cc}{\mathrm{C}}      % circle
\newcommand{\Dd}{\mathrm{B}}      % closed disc
\newcommand{\chk}[1]{\textsuperscript{[#1]}}
\definecolor{RED}{rgb}{0.8,0,0}\definecolor{red}{rgb}{0.8,0,0}
\newcommand{\placeholder}[1]{\textcolor{red}{\textbf{[#1]}}}

\title[Second-largest and minimum distance multiplicities]{Second-largest and minimum distance multiplicities in planar point sets: an upper bound of $15/11$, and related results on Erd\H{o}s Problem \#132}

\author{Wingate Jones}

\date{Draft of 24 September 2026}

\begin{document}

\begin{abstract}
For a finite set $X\subset\R^2$ let $\Dt$ and $\delta$ be the second-largest and the smallest distance determined by $X$, and let $\mult(d)$ be the number of pairs of $X$ at distance $d$. Clemen, Dumitrescu and Liu asked for
$L=\limsup_{n\to\infty}\max_{|X|=n}\min\{\mult(\Dt),\mult(\delta)\}/n$; the known bounds were $9/7\le L\le 3/2$, the upper bound coming from Vesztergombi's inequality $\mult(\Dt)\le 3n/2$. We prove that $\min\{\mult(\Dt),\mult(\delta)\}\le \tfrac{15}{11}n+C_0$ for an absolute constant $C_0$, so $L\le 15/11<3/2$. In one of the regimes of the proof (``Region~II'', defined by $\Delta\le1.94\,\Dt$ and $\tau(X)\ge\sqrt3/2+0.03$, where $\tau(X)$ is an explicit scale-invariant measure of the excess of $\Delta$ over $\Dt$ in units of $\delta$) the bound improves to $\tfrac43n+C_0$. The bound $\tfrac{15}{11}n+C_0$ is formalised in Lean~4 with Mathlib, using only the standard axioms; the refinement $\tfrac43$ is not.

We also give three results on Erd\H{o}s Problem \#132, which asks whether some distance other than the diameter must occur at most $n$ times. First, if no vertex of the convex hull has four points of $X$ at distance $\Dt$ (for instance, if no four points of $X$ are concyclic), then $\mult(\Dt)\le\tfrac32|L_1|+|L_2|$, where $L_1,L_2$ are the first two convex layers; so, under this hypothesis, the answer is yes whenever at least half as many points lie at depth $\ge 3$ as on the hull. This rests on a layer bound whose inputs from Vesztergombi's structure theory we prove directly. Second, for $n$ points in convex position there are at least three distances occurring between $1$ and $n$ times when $n$ is even and at least $6$, and when $6\le n\le18$. The cases $n=11,13,15,17$ are computer-assisted: they reduce to the unsatisfiability of finite ordinal problems, certified by DRAT (and for $n\le13$ LRAT) proof checkers, and for $n=7,11,13$ also by a Lean~4 formalisation. Third, if all but $w$ of the $n$ points in convex position lie on one circle and $n+w\ge91$, $n\ge3w$, then at least $(n-3w+1)/4$ distances are rare; the proof combines Kneser's theorem with a bound of Beukers and Smyth on torsion points of plane curves.
\end{abstract}

\maketitle

\section*{AI disclosure}
The author used an AI assistant (Claude, Anthropic) for exploration, computation, formalisation and drafting. The checks performed on the results are described in \cref{app:verification}.

\section{Introduction}\label{sec:intro}

Let $X\subset\R^2$ be a set of $n$ points. Write $\Delta=\Delta(X)$ for its diameter, $\Dt$ for its second-largest distance and $\delta$ for its minimum distance, and write $\mult(d)=\mult_X(d)$ for the number of unordered pairs of points of $X$ at distance $d$. Hopf and Pannwitz~\cite{HP34} proved that $\mult(\Delta)\le n$.

\subsection*{Erd\H{o}s Problem \#132}
Erd\H{o}s and Pach asked whether there must be a \emph{second} distance, besides the diameter, that occurs at least once and at most $n$ times, and whether the number of such distances must tend to infinity (perhaps being at least $n^{1-o(1)}$). This is Problem \#132 on the Erd\H{o}s problems website~\cite{EP132}, where it is recorded as open and where Erd\H{o}s's offer of \$100 for ``any nontrivial result'' is recorded (reference [Er97e] there). The statement fails for $n=4$ (two equilateral triangles glued along an edge) and holds for $n=5,6$ by Erd\H{o}s and Fishburn~\cite{EF95a}. We call a distance $d$ \emph{rare} if $1\le\mult(d)\le n$.

Clemen, Dumitrescu and Liu~\cite{CDL25} proved that the answer to the first question is yes for sets in convex position (their Theorem~1.2), and that $\mult(\Dt)\le n$ whenever
\begin{equation}\label{eq:cdl13}
\min\Bigl\{\tfrac32(|L_1|+|L_2|),\ \tfrac43|L_1|+2|L_2|,\ 2|L_1|+|L_2|\Bigr\}\le n,
\end{equation}
where $L_1,L_2$ are the first two convex layers (their Theorem~1.3). They showed that neither $\Dt$ nor $\delta$ is always rare: they constructed sets with $\min\{\mult(\Dt),\mult(\delta)\}\ge 9n/8-o(n)$ (their Proposition~1.5). They asked (their Problem~1.6) for the value of
\[
L:=\limsup_{n\to\infty}\ \max_{|X|=n}\ \frac{\min\{\mult(\Dt),\mult(\delta)\}}{n}.
\]
Vesztergombi~\cite{Ve87} proved that $\mult(\Dt)\le\frac32 n$ for every $n$-point planar set, which gives $L\le 3/2$. On the Erd\H{o}s problems forum, the user ienjoymath~\cite{ienjoymath} gave a construction showing $L\ge 9/7$ and conjectured $L=9/7$. As far as we know, $9/7\le L\le 3/2$ were the best bounds before this paper.

\subsection*{Results}
Our main result is an upper bound strictly below $3/2$. For a set with at least two distances put
\[
\tau(X):=\frac{\Delta^2-\Dt^2-\delta^2}{2\,\Dt\,\delta},
\]
a scale-invariant measure of how much the diameter exceeds $\Dt$ (in the scaling $\delta=1$ used in the proof, $\tau=t-(1-t^2)/(2\Dt)$ with $t=\Delta-\Dt$; see \cref{sec:setup}).

\begin{mainthm}
There is an absolute constant $C_0$ such that every $n$-point set $X\subset\R^2$ satisfies
\begin{enumerate}[label=\textup{(\alph*)}]
\item $\min\{\mult(\Dt),\mult(\delta)\}\ \le\ \tfrac{15}{11}\,n+C_0$;
\item if moreover $\Delta\le1.94\,\Dt$ and $\tau(X)\ge\sqrt3/2+3/100$, then $\min\{\mult(\Dt),\mult(\delta)\}\le\tfrac43\,n+C_0$.
\end{enumerate}
Consequently $L\le 15/11\approx 1.3636$.
\end{mainthm}
(Here $X$ is assumed to determine at least two distances; otherwise $n\le3$.)

The constant $C_0$ is explicit; the proof gives $C_0=3N_0+2\cdot10^8$ with $N_0=1.51\cdot 10^9$ (\cref{sec:assembly}). We have made no attempt to optimise it. Part~(a) is formalised in Lean~4 with Mathlib: the formal proof uses only the standard axioms (no \texttt{sorry}, no \texttt{native\_decide}) and is checked by the Lean kernel and replayed by the independent checker \texttt{leanchecker}. In three places it takes a different route from the text (the $R$-term bound, the case $\tau\le\sqrt3/2$ of Region~III, and \cref{lem:C}); see \cref{sec:sat}. Part~(b) is not formalised. A weaker bound, $\frac{54}{37}n+C_0$ with $C_0=3N_0+10^8$, obtained by a different argument in Regions II and III and with the trivial bound $\mult(\delta)\le e(S)+6|R|$ in place of \cref{sec:Rpoints}, was formalised first; that argument is written out in \cref{app:lean54}.

\Cref{thm:layer} and \cref{cor:A1} give a new class of sets for which the answer to the first question of Problem \#132 is yes. Let $m_3:=n-|L_1|-|L_2|$ be the number of points of depth at least $3$. Condition~\eqref{eq:cdl13} contains the case $2|L_1|+|L_2|\le n$, i.e.\ $|L_1|\le m_3$. We show that $|L_1|\le 2m_3$ suffices under a genericity hypothesis that holds, for instance, when no four points of $X$ are concyclic. The inputs we need from Vesztergombi's structure theory~\cite{Ve96} are proved here from scratch (\cref{prop:core2,thm:deg4}).

\Cref{sec:convex} treats the second question of Problem \#132 in convex position. \Cref{thm:convex-summary} states that a set of $n$ points in convex position has at least three rare distances for every $n$ with $6\le n\le 18$. Some of these cases follow from the literature, some follow from a parity-sharp version of the counting argument of~\cite{CDL25}, and the cases $n=11,13,15,17$ are new and computer-assisted. \Cref{thm:AC} shows that a set in convex position with all but $w$ points on one circle has at least $(n-3w+1)/4$ rare distances once $n+w\ge91$ (and $n\ge3w$); so the second question of Problem \#132 has a positive answer for such sets when $w=o(n)$. In \cref{sec:open} we list open problems.

\subsection*{Idea of the proof of Theorem 1}
Let $S$ be the set of points that have a partner at distance $\Dt$, let $Q\subseteq S$ be the set of points with exactly one such partner, and put $s=|S|$, $p=|Q|$, $r=n-s$. Vesztergombi's bound applied to $S\setminus Q$ gives $\mult(\Dt)\le\frac32s-\frac12p$. On the $\delta$ side we use the exact identity $\mult(\delta)=e(S)+\sum_{y\in R}\frac12(\deg y+m_y)$, where $e(S)$ is the number of $\delta$-pairs inside $S$, $R:=X\setminus S$, $\deg y$ is the number of $\delta$-neighbours of $y$ and $m_y$ the number of those in $S$. We prove two estimates: $e(S)\le s+p+O(1)$, and $\deg y+m_y\le10$ for all but $O(1)$ points $y\in R$ (and $\le8$ in Region~II). The combination with weights $\frac8{11}$ and $\frac3{11}$ gives $\frac{15}{11}n+O(1)$ (\cref{lem:lpnew}).

For the first estimate: almost all points of $S$ lie on the boundary of a convex body, either the ``ball polygon'' $K=\bigcap_{w\in X}\Dd(w,\Dt)$ or the convex hull, and each has a well-defined inward direction. All but $O(1)$ of the $\delta$-edges inside $S$ join points with almost equal inward directions. Such edges are either tangential edges, of which there are at most $s$, or \emph{rungs} joining the two boundaries, whose geometry is fixed by an exact identity (\cref{lem:identities}(c)). A corner lemma shows that the inner end of a rung has only one $\Dt$-partner, up to $O(1)$ exceptions, unless a second partner sits in one exactly determined reflected position; in that case the rung forces a missing tangential edge elsewhere. So every rung is paid for, either on the $\Dt$ side (through $p$) or by a missing tangential edge. For the second estimate: a point of $R$ that is not a diametral endpoint lies in the interior of $K$, while its neighbours in $S$ do not; where $\partial K$ is nearly flat on a long stretch, this confines those neighbours to an open arc of less than $240^\circ$, so there are at most $4$ of them, and the points $y$ near which $\partial K$ is not flat are $O(1)$ in number. The regime where the diameter is much larger than $\Dt$ is handled separately by an incidence argument.

\medskip\noindent\textbf{Notation.} $\Dd(c,r)$ and $\Cc(c,r)$ are the closed disc and the circle with centre $c$ and radius $r$. For unit vectors $u,u'$, $\angle(u,u')\in[0,\pi]$ is the angle between them; recall $|u-u'|\le\angle(u,u')$. Angles are in radians unless marked with a degree sign.

\section{Set-up for the proof of Theorem 1}\label{sec:setup}

We scale so that $\delta=1$. Throughout \cref{sec:setup,sec:reductions,sec:bad,sec:regions,sec:Rpoints,sec:degenerate,sec:assembly} and \cref{app:lean54}, $X$ is an $n$-point set with $n\ge 4$ and at least two distinct distances, so $\Delta>\Dt\ge\delta=1$. We put
\[
t:=\Delta-\Dt>0,\qquad \tau:=t-\frac{1-t^2}{2\Dt},\qquad \beta:=\frac1{100}.
\]
Then $\tau=(\Delta^2-\Dt^2-1)/(2\Dt)$, which is the quantity $\tau(X)$ of the introduction, and $\tau-t=(t^2-1)/(2\Dt)$. Hence $\tau\le t$ when $t\le 1$, and $\tau>1$ if and only if $t>1$.\chk{C2} A crucial feature of $X$ is that it has \emph{no distance in the open interval $(\Dt,\Delta)$}.

We use the following objects.
\begin{itemize}[leftmargin=2em]
\item $D$ is the set of endpoints of diametral pairs. Each point of $D$ is a vertex of $P:=\conv X$.
\item $S$ is the set of points of $X$ that have a partner at distance $\Dt$; $R:=X\setminus S$ and $r:=|R|$, $s:=|S|$.
\item $Q\subseteq S$ is the set of points with exactly one partner at distance $\Dt$ (\emph{$\Dt$-degree} $1$), and $p:=|Q|$.
\item $K:=\bigcap_{w\in X}\Dd(w,\Dt)$, a compact convex set (a ``ball polygon'').
\item $e(Y)$ is the number of $\delta$-pairs inside $Y\subseteq X$. For $y\in X$, $\deg y$ is the number of $\delta$-neighbours of $y$, and for $y\in R$, $m_y$ is the number of $\delta$-neighbours of $y$ that lie in $S$.
\end{itemize}
The \emph{degenerate regime} is $\Delta>1.94\,\Dt$ (\cref{sec:degenerate}); the \emph{nondegenerate regime} is $\Delta\le 1.94\,\Dt$ (\cref{sec:reductions,sec:bad,sec:regions,sec:Rpoints}).

\begin{lemma}[ball polygon]\label{lem:ballpoly}
\begin{enumerate}[label=(\alph*)]
\item $X\setminus D\subseteq K$.
\item $S\setminus D\subseteq\partial K$.
\item $\dist(z,K)\ge t$ for every $z\in D$.
\end{enumerate}
\end{lemma}
\begin{proof}
(a) A point $y\notin D$ has all its distances at most $\Dt$, since there are none in $(\Dt,\Delta)$. (b) If $v\in S\setminus D$ and $|vw|=\Dt$, then $v\in K\cap \Cc(w,\Dt)\subseteq\partial K$. (c) If $z\in D$ and $|zx|=\Delta$, then $K\subseteq\Dd(x,\Dt)$, while $|zx|=\Dt+t$.
\end{proof}

\subsection*{Direction field}
For each $p\in S$ we fix a unit vector $u_p$.
\begin{itemize}[leftmargin=2em]
\item If $p\in S\setminus D$, choose a $\Dt$-partner $w_p$ and put $u_p:=(w_p-p)/\Dt$. Since $K\subseteq\Dd(w_p,\Dt)$, $u_p$ is an inward normal of $K$ at $p$.
\item If $p\in S\cap D$, choose a diametral partner $x_p$ and put $u_p:=(x_p-p)/\Delta$. This is an inward normal of $P$ at $p$.
\end{itemize}

\begin{lemma}[exact identities]\label{lem:identities}
Let $y\in X$, $y\neq p$.
\begin{enumerate}[label=(\alph*)]
\item If $p\in S\cap D$, then $(y-p)\cdot u_p\ge |y-p|^2/(2\Delta)>0$.
\item If $p\in S\setminus D$, then either $(y-p)\cdot u_p\ge |y-p|^2/(2\Dt)$ (we say $y$ is \emph{normal} for $p$), or $|y-w_p|=\Delta$, so that $y\in D$ (we say $y$ is an \emph{exception} for $p$).
\item In the exception case with $|y-p|=1$ we have exactly $(y-p)\cdot u_p=-\tau$. Hence $y-p=-\tau u_p\pm\sigma u_p^{\perp}$ with $\sigma:=\sqrt{1-\tau^2}$. In particular exceptions at distance $1$ exist only when $\tau\le 1$.
\end{enumerate}
\end{lemma}
\begin{proof}
(a) $|y-x_p|^2=|y-p|^2-2\Delta(y-p)\cdot u_p+\Delta^2\le\Delta^2$. (b) $|y-w_p|^2=|y-p|^2-2\Dt(y-p)\cdot u_p+\Dt^2$, and $|y-w_p|\in[0,\Dt]\cup\{\Delta\}$. (c) With $|y-p|=1$ and $|y-w_p|=\Delta=\Dt+t$ the same identity gives $(y-p)\cdot u_p=(1+\Dt^2-\Delta^2)/(2\Dt)=(1-2\Dt t-t^2)/(2\Dt)=-\tau$.\chk{C2}
\end{proof}
The proof uses only $|p-x_p|=\Delta$, resp.\ $|p-w_p|=\Dt$. So (a) holds for every $p\in D$ (also $p\in R\cap D$) and every diametral partner $x$ of $p$ in place of $x_p$, and (b), (c) hold for every $\Dt$-partner $w$ of $p$ in place of $w_p$, with $u_p$ replaced by $(w-p)/\Dt$. We use these variants below.

\begin{lemma}[no common normal]\label{lem:nocommon}
If $p\neq q$ both lie in $S\setminus D$, then $u_p\neq u_q$.
\end{lemma}
\begin{proof}
Neither point lies in $D$, so each is normal for the other: $(q-p)\cdot u_p>0$ and $(p-q)\cdot u_q>0$. These cannot both hold if $u_p=u_q$.
\end{proof}

\subsection*{Good edges, T-edges and rungs}
A $\delta$-edge $pq$ inside $S$ is \emph{good} if $\angle(u_p,u_q)\le\beta$, and \emph{bad} otherwise. For a good edge put $e:=q-p$, $a:=e\cdot u_p$ and $b:=-e\cdot u_q$. Then $|a+b|=|e\cdot(u_p-u_q)|\le\beta$.

\begin{lemma}[dichotomy]\label{lem:dichotomy}
Assume $\tau>2\beta$. Each good edge is of exactly one of the following two kinds.
\begin{itemize}[leftmargin=2em]
\item \emph{T-edge:} $a,b\in(0,\beta)$.
\item \emph{Rung:} one end $v\in S\setminus D$ sees the other end $z\in S\cap D$ as an exception. Then $(z-v)\cdot u_v=-\tau$ and $(v-z)\cdot u_z\in[\tau-\beta,\tau+\beta]$.
\end{itemize}
\end{lemma}
\begin{proof}
By \cref{lem:identities}, from each end the other end is either normal (then the corresponding one of $a,b$ is positive) or an exception (then it equals $-\tau$). An exception target lies in $D$, and points of $D$ see everything as normal, so both ends cannot be exceptions. If neither is, then $a,b>0$ and $a+b\le\beta$. If one is, say $a=-\tau$, then $b\in[\tau-\beta,\tau+\beta]$, and $b>0$ shows that $b$ is not an exception value.
\end{proof}

We say a T-neighbour $q$ of $p$ is on the \emph{right} (resp.\ \emph{left}) of $p$ if $(q-p)\cdot u_p^{\perp}>0$ (resp.\ $<0$), where $u^\perp$ denotes $u$ rotated by $-90^\circ$.
\begin{itemize}[leftmargin=2em]
\item[(T1)] Each $p$ has at most one T-neighbour on each side. Indeed, two T-directions on one side both make an angle in $(0,\arcsin\beta)$ with $\pm u_p^\perp$, so they are less than $60^\circ$ apart, whereas two points at distance $1$ from $p$ and at distance $\ge 1$ from each other are at least $60^\circ$ apart as seen from $p$.
\item[(T2)] If $q$ is the right T-neighbour of $p$ then $p$ is the left T-neighbour of $q$, since $\angle(u_p,u_q)\le\beta$.
\end{itemize}
Hence the T-edges form vertex-disjoint directed paths and cycles, and every cycle has length at least $3$ (a vertex cannot be both the left and the right T-neighbour of another). We say each point of $S$ has a \emph{left slot} and a \emph{right slot}; a T-edge from $p$ to its right T-neighbour $q$ uses the right slot of $p$ and the left slot of $q$. So the number $e_T$ of T-edges is half the number of used slots, and $e_T\le s$. T-edges with both ends in $S\setminus D$ or both in $S\cap D$ are \emph{pure}; the others are \emph{mixed}. Note that (T1) and (T2) apply to mixed T-edges too.

\subsection*{Turning}
Let $\gamma$ be a closed arc of the boundary of a compact convex set with nonempty interior. Its \emph{turning} $\turn(\gamma)$ is the length of the set $N(\gamma)\subseteq S^1$, the union of the normal cones at \emph{all} points of $\gamma$, including the full normal cones at its endpoints.

\begin{lemma}[multiplicity]\label{lem:multiplicity}
Let $\gamma_1,\dots,\gamma_m$ be closed boundary arcs of one convex body, and suppose every boundary point lies in at most $M$ of them. Then $\sum_i\turn(\gamma_i)\le 2\pi M$.
\end{lemma}
\begin{proof}
For all but countably many directions $\nu\in S^1$ the set of boundary points with outer (or inner) normal $\nu$ is a single point $y_\nu$, and $\nu\in N(\gamma_i)$ if and only if $y_\nu\in\gamma_i$. Integrate over $\nu$.
\end{proof}

Points of $S\setminus D$ lie on $\partial K$ and points of $S\cap D$ lie on $\partial P$; in each case distinct points are at Euclidean, hence arc-length, distance at least $1$. So a boundary arc of length $\ell$ contains at most $\ell+1$ of them, and a boundary of perimeter $\ell$ contains at most $\ell$ of them.

\section{Reductions}\label{sec:reductions}

\begin{lemma}[(VS)]\label{lem:VS}
For every $S'\subseteq S$, the number of $\Dt$-pairs inside $S'$ is at most $\frac32|S'|$. In particular $\mult(\Dt)\le\frac32 s$.
\end{lemma}
\begin{proof}
Every $\Dt$-pair of $X$ lies in $S$, and every distance within $S'$ is a distance of $X$. If $S'$ contains no $\Dt$-pair there is nothing to prove. If $S'$ contains a diametral pair of $X$, then $\operatorname{diam}S'=\Delta$ and, because $X$ has no distances in $(\Dt,\Delta)$, $\Dt$ is the second-largest distance of $S'$; Vesztergombi's theorem~\cite{Ve87} gives at most $\frac32|S'|$ such pairs. Otherwise every distance in $S'$ is at most $\Dt$, so $\Dt=\operatorname{diam}S'$ and Hopf--Pannwitz~\cite{HP34} gives at most $|S'|$.
\end{proof}

We use Vesztergombi's theorem in exactly the form of its main theorem: for any set of $m$ points in the plane, the second-largest distance occurs at most $\frac32 m$ times~\cite[Theorem]{Ve87}. There is no threshold on $m$ and no additive constant.

\begin{lemma}[weighted $\Dt$ bound]\label{lem:VSQ}
$\mult(\Dt)\le p+\frac32(s-p)=\frac32s-\frac12p$.
\end{lemma}
\begin{proof}
At most $p$ $\Dt$-pairs touch $Q$, since each point of $Q$ lies in exactly one $\Dt$-pair. Every other $\Dt$-pair lies inside $S\setminus Q$, and there are at most $\frac32(s-p)$ of them by \cref{lem:VS}.
\end{proof}

\begin{lemma}[exact identity]\label{lem:identity}
$\mult(\delta)=e(S)+\sum_{y\in R}\tfrac12(\deg y+m_y)$. Moreover $\deg y\le 6$ and $m_y\le\deg y$, so $\deg y+m_y\le12$ for every $y\in R$.
\end{lemma}
\begin{proof}
The $\delta$-pairs inside $S$ number $e(S)$. The pairs with one end in $S$ and one in $R$ number $\sum_{y\in R}m_y$. The pairs inside $R$ are counted twice in $\sum_{y\in R}(\deg y-m_y)$. Adding, $\mult(\delta)=e(S)+\sum_ym_y+\frac12\sum_y(\deg y-m_y)$, which is the identity. Each point has at most $6$ points at distance $\delta$, since these are pairwise at least $\delta$ apart.
\end{proof}

\begin{corollary}[(E)]\label{lem:E}
$\mult(\delta)\le e(S)+6r$. \qed
\end{corollary}

\begin{lemma}[linear combination]\label{lem:lpnew}
Suppose $e(S)\le s+p+C$, and that $\deg y+m_y\le10$ for all but at most $C'$ points $y\in R$. Then
\[
\min\{\mult(\Dt),\mult(\delta)\}\le\tfrac{15}{11}\,n+C+C'.
\]
If instead $\deg y+m_y\le 8$ for all but at most $C'$ points $y\in R$, then $\min\{\mult(\Dt),\mult(\delta)\}\le\frac43n+C+C'$.
\end{lemma}
\begin{proof}
For an exceptional $y$ we still have $\deg y+m_y\le12$ (\cref{lem:identity}). So in the first case $\mult(\delta)\le s+p+5r+C+C'$, and with \cref{lem:VSQ}
\[
\min\{\mult(\Dt),\mult(\delta)\}\le\tfrac8{11}\bigl(\tfrac32s-\tfrac12p\bigr)+\tfrac3{11}\bigl(s+p+5r+C+C'\bigr)=\tfrac{15}{11}(s+r)-\tfrac1{11}p+\tfrac3{11}(C+C').
\]
In the second case $\mult(\delta)\le s+p+4r+C+2C'$, and the weights $\frac23,\frac13$ give $\frac43(s+r)+\frac13(C+2C')$.\chk{LP}
\end{proof}

\begin{remark}\label{rem:lpvalues}
If $\deg y+m_y\le 2c$ for all but $O(1)$ points of $R$, the best weights give $\max\{\frac43,\frac{3c}{2c+1}\}\,n+O(1)$ under the hypothesis $e(S)\le s+p+O(1)$, and $\frac{3c}{2c+1}\,n+O(1)$ under the stronger hypothesis $e(S)\le s+O(1)$; the weight on $\mult(\Dt)$ must be at least $\frac23$ to cancel the $p$-term.\chk{LP} For $2c=12,10,8,6$ these values are $\frac{18}{13},\frac{15}{11},\frac43,\frac43$, respectively $\frac{18}{13},\frac{15}{11},\frac43,\frac97$. In particular the trivial bound (E) together with $e(S)\le s+p+O(1)$ already gives $\frac{18}{13}n+O(1)$, with weights $\frac{10}{13},\frac3{13}$.
\end{remark}

So in the nondegenerate regime it suffices to show $e(S)\le s+p+O(1)$ in every $\tau$-regime (\cref{sec:regions}), and $\deg y+m_y\le 10$ (resp.\ $\le 8$ in Region~II) for all but $O(1)$ points $y\in R$ (\cref{sec:Rpoints}).

\subsection*{Small cases}
Suppose we are in the nondegenerate regime and $\Dt<10^4$. Then $\Delta<1.94\cdot10^4$, so $X$ lies in a disc of radius $\Delta$ about any of its points, and since the discs of radius $\frac12$ about points of $X$ are disjoint, $n\le(2\Delta+1)^2<N_0:=1.51\cdot10^9$.\chk{C11} Since $\mult(\delta)\le 3n$, Theorem~1 holds for such $X$ once $C_0\ge 3N_0$. \emph{From now on, in the nondegenerate regime, $\Dt\ge 10^4$}; in particular $2/\Dt\le\beta/2$.\chk{C3}
\section{Bad edges are \texorpdfstring{$O(1)$}{O(1)} in the nondegenerate regime}\label{sec:bad}

\subsection{A lens lemma}

\begin{lemma}[lens]\label{lem:lens}
Let $\Lambda=\Dd(c_1,\rho)\cap\Dd(c_2,\rho)$ with $0<|c_1c_2|=d<2\rho$, and put $\cos\varphi_0:=d/(2\rho)$. Assume
\begin{equation}\label{eq:lenshyp}
\rho\ \ge\ \frac{2}{\cos(\varphi_0/2)-\cos\varphi_0}.
\end{equation}
Let $p$ lie on the arc of $\partial\Lambda$ belonging to $\Cc(c_1,\rho)$, let $q$ lie on the arc belonging to $\Cc(c_2,\rho)$, and suppose $|pq|\le 1$. Then there is a corner $c$ of $\Lambda$ such that the boundary arcs from $p$ to $c$ and from $c$ to $q$ each have length at most $\rho_0:=1/\sin(\varphi_0/2)$.
\end{lemma}
\begin{proof}
Take $c_1=(-d/2,0)$, $c_2=(d/2,0)$, $p=c_1+\rho(\cos\varphi,\sin\varphi)$ and $q=c_2+\rho(-\cos\psi,\sin\psi)$ with $|\varphi|,|\psi|\le\varphi_0$.
\emph{Step 1.} $p_x-q_x=\rho(\cos\varphi-\cos\varphi_0)+\rho(\cos\psi-\cos\varphi_0)$ with both terms $\ge0$; since $|p_x-q_x|\le1$, each term is at most $1$. By~\eqref{eq:lenshyp}, $\cos\varphi<\cos(\varphi_0/2)$, i.e.\ $|\varphi|>\varphi_0/2$; likewise $|\psi|>\varphi_0/2$.
\emph{Step 2.} By the mean value theorem, $\rho(\varphi_0-|\varphi|)=\rho(\cos\varphi-\cos\varphi_0)/\sin\xi\le 1/\sin\xi$ for some $\xi\in(\varphi_0/2,\varphi_0)\subseteq(0,\pi/2)$, so $\rho(\varphi_0-|\varphi|)\le\rho_0$; likewise for $\psi$.
\emph{Step 3.} If $\varphi,\psi$ had opposite signs, then $|p_y-q_y|\ge 2\rho\sin(\varphi_0/2)$. Now~\eqref{eq:lenshyp} gives $\rho\ge 2/(1-\cos\varphi_0)=1/\sin^2(\varphi_0/2)$, so $2\rho\sin(\varphi_0/2)\ge 2/\sin(\varphi_0/2)\ge2>1$, a contradiction.
\emph{Step 4.} So the corner $c=(0,\operatorname{sign}(\varphi)\rho\sin\varphi_0)$ works.
\end{proof}

\begin{lemma}[lens-corner arc length]\label{lem:lenscorner}
In the situation of \cref{lem:lens}, suppose moreover that $p,q\in\partial K'$ for a compact convex $K'\subseteq\Lambda$, and $|pq|=1$. Let $H$ be the closed half-plane bounded by the line $pq$ that contains $c$, and let $\gamma:=\partial K'\cap H$ be the boundary arc of $K'$ from $p$ to $q$ on $c$'s side. Then $\operatorname{length}(\gamma)\le 2\rho_0$.
\end{lemma}
\begin{proof}
$\partial\Lambda$ contains no segment, so the line $pq$ meets it only in $p$ and $q$, and $\partial\Lambda\cap H$ is exactly the arc $p\to c\to q$, of length at most $2\rho_0$. Since $K'\cap H\subseteq\Lambda\cap H$ are convex, monotonicity of perimeter gives $\operatorname{length}(\gamma)+1\le 2\rho_0+1$. (If $K'\cap H$ is the segment $[p,q]$, then $\gamma$ is that segment, of length $1\le2\rho_0$.)
\end{proof}

\emph{Constants.}\chk{C7} For \emph{KK lenses} ($\rho=\Dt$, $d\le\Delta\le1.94\Dt$) we have $\cos\varphi_0\le 0.97$, so $\rho_0^2=2/(1-\cos\varphi_0)\le 66.7<8.2^2$. The function $c\mapsto\cos(\varphi_0/2)-c=\sqrt{(1+c)/2}-c$ is decreasing on $[0,1)$, so~\eqref{eq:lenshyp} holds as soon as $\rho\ge 89.01$, in particular for $\rho=\Dt\ge10^4$. For \emph{DD lenses} ($\rho=\Delta$, $d\le\Delta$) we have $\cos\varphi_0\le\frac12$, $\rho_0\le 2$, and~\eqref{eq:lenshyp} needs only $\rho\ge 5.5$.

\subsection{The count}
We split bad edges $pq$ by type: \emph{KK} (both ends in $S\setminus D$), \emph{DD} (both in $S\cap D$) and \emph{cross} (one end in each).

If the perimeter of $K$ is below $1202$, then $|S\setminus D|<1202$ and at most $6\cdot1202$ edges touch $S\setminus D$; similarly for $P$ and $S\cap D$. We may therefore assume both perimeters are at least $1202$ when counting the bad edges touching the corresponding row; these small cases contribute a constant below $1.5\cdot10^4$ in total.

\subsubsection*{KK edges}
Let $p,q\in S\setminus D$ form a bad edge. If $w_p=w_q$, then $\angle(u_p,u_q)\le2\arcsin(1/(2\Dt))<\beta$,\chk{C3} so the edge would be good; hence $w_p\neq w_q$. Then $K\subseteq\Lambda:=\Dd(w_p,\Dt)\cap\Dd(w_q,\Dt)$, $p$ lies on the $\Cc(w_p,\Dt)$-arc and $q$ on the $\Cc(w_q,\Dt)$-arc of $\partial\Lambda$, and $d=|w_pw_q|\le\Delta\le1.94\Dt$. \Cref{lem:lens,lem:lenscorner} give an arc $\gamma_{pq}\subseteq\partial K$ from $p$ to $q$ of length at most $2\rho_0$. The union of the normal cones along the closed arc $\gamma_{pq}$ is connected and contains $u_p$ and $u_q$, so $\turn(\gamma_{pq})\ge\angle(u_p,u_q)>\beta$. If $y\in\gamma_{pq}$ then $p$ lies within arc distance $2\rho_0$ of $y$; there are at most $4\rho_0+1$ such $p$, each of degree at most $6$, so every point lies on at most $M=6(4\rho_0+1)$ of the arcs. By \cref{lem:multiplicity},
\[
\#\text{KK-bad}\le 12\pi(4\rho_0+1)/\beta<1.28\cdot10^5.\chk{C9}
\]

\subsubsection*{DD edges}
The same argument, on $\partial P$ with the lens $\Dd(x_p,\Delta)\cap\Dd(x_q,\Delta)$ and $\rho_0\le 2$, gives $\#\text{DD-bad}\le 12\pi\cdot 9/\beta<3.4\cdot10^4$. (If $x_p=x_q$ the angle is less than $2/\Delta<\beta$ and the edge is good.)

\subsubsection*{Cross edges}
Let $zv$ be a bad edge with $z\in S\cap D$ and $v\in S\setminus D$. Since $z\in X$, $|zw_v|\in[0,\Dt]\cup\{\Delta\}$.

\emph{(c1) $|zw_v|=\Delta$.} By \cref{lem:identities}(a), $(w_v-z)/\Delta$ is a second inward normal of $P$ at $z$. Its angle to $u_v$ is the angle at $w_v$ of the triangle $zvw_v$; that angle is opposite the side of length $1$, hence acute, and by the law of sines it is at most $\arcsin(1/\Delta)<2/\Dt\le\beta/2$.\chk{C3} As $\angle(u_z,u_v)>\beta$, the normal cone of $P$ at $z$ has angle greater than $\beta/2$. Normal cones at distinct vertices have disjoint interiors, so fewer than $4\pi/\beta$ points $z$ qualify, each of degree at most $6$: $\#(\mathrm{c1})<24\pi/\beta<7600$.

\emph{(c2) $|zw_v|\le\Dt$.} Let $g$ be the unit vector from $v$ towards $x_z$, and $\alpha:=\angle(u_v,g)$.
\begin{enumerate}[leftmargin=2em,label=(\roman*)]
\item \emph{Depth.} $v\notin D$ gives $|vx_z|\le\Dt$, while $|zx_z|=\Delta$ and $|zv|=1$ give $|vx_z|\ge\Delta-1$. So $v\in\Dd(x_z,\Dt)$ at depth $d_0:=\Dt-|vx_z|\in[0,1-t]$; in particular this case needs $t\le 1$. Also $K\subseteq\Dd(x_z,\Dt)$.
\item \emph{$\alpha$ is not small.} $\angle(u_z,g)$ is the angle at $x_z$ of the triangle $zvx_z$, which is at most $\arcsin(1/\Delta)<\beta/2$; hence $\alpha\ge\angle(u_v,u_z)-\angle(u_z,g)>\beta/2$.
\item \emph{$\alpha$ is not close to $\pi$.} Since $w_v,x_z\in X$, $|w_vx_z|\le\Delta\le1.94\Dt$. The law of cosines in the triangle $vw_vx_z$ (sides $\Dt$ and $\Dt-d_0$ at $v$, enclosing the angle $\alpha$) gives $\cos\alpha\ge(1+\varrho^2-1.94^2)/(2\varrho)$, where $\varrho:=(\Dt-d_0)/\Dt\in[1-10^{-4},1]$. The right side is increasing in $\varrho$, so $\cos\alpha\ge-0.8821$, i.e.\ $\alpha\le151.9^\circ$. By concavity of $\sin$ on $[0,\pi]$, $\sin\alpha\ge\min\{\sin(\beta/2),0.471\}=\sin(\beta/2)$.\chk{C8}
\item \emph{The walk.} Let $e_0\perp u_v$ be the unit vector with $e_0\cdot g=-\sin\alpha<0$. Follow $\partial K$ from $v$ in the direction whose forward one-sided tangent is closest to $e_0$, parametrised by arc length, $\gamma(0)=v$. Suppose the turning of $\gamma([0,\sigma])$ is at most $\beta/4$. This turning includes the whole normal cone $N_K(v)$. Then every normal along the arc is within $\beta/4$ of $u_v$, so every forward tangent $\gamma'$ is within $\beta/4$ of $e_0$ (at a corner $v$, the forward tangent is perpendicular to an extreme normal of $N_K(v)$, which is also within $\beta/4$ of $u_v$). Hence $\gamma'\cdot g\le-\sin\alpha+2\sin(\beta/8)\le-\sin(\beta/2)+2\sin(\beta/8)<-\beta/5$.\chk{C8} Integrating, $(\gamma(\sigma)-v)\cdot g\le-\sigma\beta/5$, and so
\[
|\gamma(\sigma)-x_z|^2\ \ge\ (\Dt-d_0)^2+\tfrac25(\Dt-d_0)\,\sigma\beta .
\]
For $\sigma=L:=6/\beta=600$ the right side exceeds $\Dt^2$,\chk{C8} which puts $\gamma(L)$ outside $\Dd(x_z,\Dt)\supseteq K$, a contradiction. Hence the closed arc $J_v$ of length $L$ from $v$ in the chosen direction has turning greater than $\beta/4$.
\item \emph{Count.} A point $y$ lies in $J_v$ only if $v$ is within arc distance $L$ of $y$, which allows at most $2L+1$ points $v$. By \cref{lem:multiplicity}, $\#\{v\}\le 2\pi(2L+1)/(\beta/4)$. Each such $v$ has degree at most $6$, so $\#(\mathrm{c2})\le48\pi(2L+1)/\beta<1.82\cdot10^7$.\chk{C9}
\end{enumerate}


\begin{proposition}\label{prop:bad}
In the nondegenerate regime the number $C_{\mathrm{bad}}$ of bad edges satisfies $C_{\mathrm{bad}}<1.83\cdot10^7$.\chk{C9}
\end{proposition}

\section{The nondegenerate regime: bounding \texorpdfstring{$e(S)$}{e(S)}}\label{sec:regions}

In this section we prove $e(S)\le s+p+C$ with an absolute constant $C$ in every $\tau$-regime of the nondegenerate regime. Rungs are counted with the bad edges excluded: $e(S)\le e_T+\rho+C_{\mathrm{bad}}$, where $\rho$ is the number of good rungs. (When $\tau\le2\beta$, \cref{lem:dichotomy} is not used.)

\subsection{Region I: \texorpdfstring{$\tau\le\sqrt3/2-2\beta$}{tau <= sqrt3/2 - 2 beta}}\label{sec:regionI}
Let $p\in S\setminus D$ and let $q$ be a good neighbour of $p$; the \emph{elevation} of $q-p$ is $\arcsin((q-p)\cdot u_p)$. If $q$ is normal for $p$, then $b>0$ (by \cref{lem:identities}(a) if $q\in D$, by (b) if $q\notin D$), so $a<\beta$. Otherwise $q$ is an exception and $a=-\tau$.
\begin{itemize}[leftmargin=2em]
\item If $\tau\ge0$, every good neighbour on one side of $\R u_p$ has elevation in $[-\arcsin\tau,\arcsin\beta)$, a window of width less than $\arcsin(\sqrt3/2-2\beta)+\arcsin\beta\approx58.35^\circ<60^\circ$.\chk{C4}
\item If $\tau<0$, then $|\tau|\le1/(2\Dt)$ and all elevations lie in $(0,\arcsin\beta)\cup\{\arcsin|\tau|\}$; the window is again less than $60^\circ$.\chk{C4}
\end{itemize}
Let $p\in S\cap D$. Then $a>0$ always, and either $a<\beta$, or $p$ is the exception of $q$ and $a\in[\tau-\beta,\tau+\beta]$; the window is at most $\arcsin(\sqrt3/2-\beta)\approx58.87^\circ<60^\circ$.\chk{C4}

Since two neighbours of $p$ are at least $60^\circ$ apart, $p$ has at most one good neighbour on each side. Hence
\[
e(S)\le s+C_{\mathrm{bad}}\qquad\text{in Region I.}
\]

\subsection{Common facts for \texorpdfstring{$\tau>\sqrt3/2-2\beta$}{tau > sqrt3/2 - 2 beta}}\label{sec:common}
Here $\tau>2\beta$,\chk{C5,C6} so \cref{lem:dichotomy} applies: every good edge is a T-edge or a rung. \emph{Region II} is $\tau\ge\sqrt3/2+3\beta$; \emph{Region III} is $\sqrt3/2-2\beta<\tau<\sqrt3/2+3\beta$. If $\tau>1$ then $t>1$, and by \cref{lem:identities}(c) there are no rungs; so $e(S)\le e_T+C_{\mathrm{bad}}\le s+C_{\mathrm{bad}}$ by (T1). \emph{From now on in this section assume $\tau\le 1$}; then $t\le1$ and $t\ge\tau$ (\cref{sec:setup}). Put
\[
\theta_0:=\arccos\tau\in[0,\pi/2),\qquad \sigma=\sqrt{1-\tau^2}=\sin\theta_0,\qquad o_\pm(u):=-\tau u\pm\sigma u^\perp .
\]
By \cref{lem:identities}(c), a rung $(v,z)$ with $v\in S\setminus D$ has $z=v+o_+(u_v)$ or $z=v+o_-(u_v)$; in particular \emph{a K-vertex has at most two rungs}. The map $u\mapsto o_\pm(u)$ is a rotation, so $|o_\pm(u)-o_\pm(u')|=|u-u'|\le\angle(u,u')$, and $|o_+(u)-o_-(u)|=2\sigma$. The \emph{K-row} is $S\setminus D$ and the \emph{D-row} is $S\cap D$; rungs join the two rows.

\emph{T-steps.} If $q$ is the right T-neighbour of $p$, then $q-p=\cos\eps\,u_p^\perp+\sin\eps\,u_p$ with $\eps\in(0,\arcsin\beta)$, since $a=(q-p)\cdot u_p\in(0,\beta)$.

\begin{lemma}[mixed T-edges]\label{lem:mixedT}
Assume $\sqrt3/2-2\beta<\tau\le1$. Let $v\in S\setminus D$ and $z\in S\cap D$ be joined by a T-edge. Then the closed arc of $\partial K$ of length $2$ starting at $v$ on $z$'s side has turning at least $t-\arcsin\beta$. Consequently there are fewer than $2\pi\cdot10/(t-\arcsin\beta)$ mixed T-edges: fewer than $71$ in Region~II and fewer than $76$ in Region~III.
\end{lemma}
\begin{proof}
Here $t\ge\tau>\sqrt3/2-2\beta>0.846$. Suppose the turning is $\Theta<t-\arcsin\beta$; note $t-\arcsin\beta<1<\pi/3$. Let $e_0$ be the unit tangent direction at $v$ perpendicular to $u_v$ pointing to $z$'s side. Every forward tangent along the arc is within $\Theta$ of $e_0$, so $|\gamma(\sigma')-v|\ge\sigma'\cos\Theta\ge\sigma'/2$ for the arc-length parametrisation $\gamma$, and some $\sigma^*\le2$ has $|\gamma(\sigma^*)-v|=1$. The unit vector $\gamma(\sigma^*)-v$ lies within $\Theta$ of $e_0$, and the unit vector $z-v$ lies within $\arcsin\beta$ of $e_0$. Hence $|\gamma(\sigma^*)-z|\le\Theta+\arcsin\beta<t\le\dist(z,K)$ by \cref{lem:ballpoly}(c). But $\gamma(\sigma^*)\in K$, a contradiction.

For the count: each $v$ has at most two T-edges, and $v$ must lie within arc distance $2$ of a point $y$ of its arc, so each boundary point lies on at most $M\le10$ of these arcs. By \cref{lem:multiplicity} the number of mixed T-edges is less than $2\pi\cdot10/(t-\arcsin\beta)$. In Region~II, $t-\arcsin\beta\ge0.886$ and the bound is below $71$;\chk{C5} in Region~III, $t-\arcsin\beta>0.836$ and the bound is below $76$.\chk{III}
\end{proof}

\begin{lemma}[corner lemma]\label{lem:corner}
Let $v\in S\setminus D$ have a rung partner $z$ and a second $\Dt$-partner $w'\neq w_v$, and put $u':=(w'-v)/\Dt$ and $c:=z-v$. Then exactly one of the following holds.
\begin{enumerate}[label=\textup{(\roman*)}]
\item $|zw'|\le\Dt$. Then $c\cdot u'\ge1/(2\Dt)>0$ and $\angle(u_v,u')>\pi/2-\theta_0$.
\item $|zw'|=\Delta$. Then $u'=u^*$, the reflection of $u_v$ in the line $\R c$, and $\angle(u_v,u')=2\theta_0>0$.
\end{enumerate}
In both cases $u_v$ and $u'$ are inward normals of $K$ at $v$, and the normal cone $N_K(v)$ has angle at least $\angle(u_v,u')$.
\end{lemma}
\begin{proof}
Since $z\in X$, $|zw'|\in[0,\Dt]\cup\{\Delta\}$, so exactly one case occurs. We have $c\cdot u_v=-\tau$ and $|c|=1$, so $\angle(c,u_v)=\pi-\theta_0$.
(i) By \cref{lem:identities}(b) applied with $w'$ in place of $w_v$, $z$ is normal for $v$ with respect to $w'$: $c\cdot u'\ge|c|^2/(2\Dt)>0$. So $\angle(c,u')<\pi/2$ and $\angle(u_v,u')\ge\angle(c,u_v)-\angle(c,u')>\pi/2-\theta_0$.
(ii) By \cref{lem:identities}(c) applied with $w'$, $c\cdot u'=-\tau$. The unit vectors $u$ with $c\cdot u=-\tau$ are the two vectors at angle $\theta_0$ from $-c$, namely $u_v$ and its reflection $u^*$ in $\R c$. Since $w'\neq w_v$ and both are at distance $\Dt$ from $v$, $u'\neq u_v$; so $u'=u^*$, $\theta_0>0$, and $\angle(u_v,u')=2\theta_0$.
Both $u_v$ and $u'$ are inward normals of $K$ at $v$, because $K\subseteq\Dd(w_v,\Dt)\cap\Dd(w',\Dt)$. Moreover $\angle(u_v,u')<\pi$, because $|w_vw'|\le\Delta<2\Dt$; so the (convex) normal cone contains the shorter arc between them.
\end{proof}
Normal cones of $K$ at distinct boundary points have disjoint interiors. We say that $v$ is \emph{of type \textup{(i)}} if some second $\Dt$-partner of $v$ is as in (i), and \emph{of type \textup{(ii)}} if $v$ has $\Dt$-degree at least $2$ and is not of type (i). A rung K-end of type (ii) has $\Dt$-degree exactly $2$, since only the directions $u_v$ and $u^*$ are available.

\subsection{Region II: \texorpdfstring{$\tau\ge\sqrt3/2+3\beta$}{tau >= sqrt3/2 + 3 beta}}\label{sec:regionII}
We keep $\tau\le1$. Then $\theta_0\le\arccos(\sqrt3/2+3\beta)<26.4^\circ<30^\circ$ and $\sigma\le\sigma_{\max}\approx0.4440$.\chk{C5}

\begin{itemize}[leftmargin=2em]
\item[(R1)] \emph{Each vertex has at most one rung.} At $v\in S\setminus D$: the two candidate positions $v+o_\pm(u_v)$ are $2\sigma<1$ apart. At $z\in S\cap D$: if $z$ had rungs to $v\neq v'$, both rungs are good, so $\angle(u_v,u_{v'})\le2\beta$ and $|v-v'|=|o_{\pm}(u_v)-o_{\pm'}(u_{v'})|\le2\sigma+2\beta\approx0.908<1$,\chk{C5} a contradiction.
\end{itemize}
So rungs form a matching. By \cref{lem:corner}, a type-(i) vertex has a normal cone of angle greater than $\pi/2-\theta_0>\pi/3$, so \emph{there are at most $5$ type-\textup{(i)} vertices}. Every other rung K-end either lies in $Q$ or is of type (ii).

\begin{lemma}[far partners persist]\label{lem:B}
Let $z\in D$ and $w\in X$ with $|zw|=\Delta$, put $\hat w:=(w-z)/\Delta$, and let $y\in X$ with $|zy|=1$ and $\eta:=(y-z)\cdot\hat w$. Then $|yw|^2=\Delta^2+1-2\Delta\eta$. If
\[
\eta<\frac{\Delta^2-\Dt^2+1}{2\Delta}=t+\frac{1-t^2}{2\Delta},
\]
then $|yw|>\Dt$, hence $|yw|=\Delta$. When $t\le1$ the threshold is at least $t$.
\end{lemma}
\begin{proof}
The identity is $|y-w|^2=|y-z|^2-2(y-z)\cdot(w-z)+|w-z|^2$. The inequality $\Delta^2+1-2\Delta\eta>\Dt^2$ is the displayed condition, and $\Delta^2-\Dt^2=2\Delta t-t^2$ gives the second form of the threshold. There are no distances in $(\Dt,\Delta)$.\chk{II}
\end{proof}

\begin{lemma}\label{lem:C}
Let $(v,z)$ be a rung whose K-end $v$ is of type \textup{(ii)}. Then $z$ has no T-neighbour at all, pure or mixed.
\end{lemma}
\begin{proof}
Let $w:=w_v$ and let $w^*$ be the second $\Dt$-partner of $v$, in direction $u^*$ from $v$. Since the rung is an exception for $v$ with respect to both partners (\cref{lem:corner}(ii)), $|zw|=|zw^*|=\Delta$. Put $\hat w:=(w-z)/\Delta$ and $\hat w^*:=(w^*-z)/\Delta$. The angle between $\hat w$ and $u_v$ is the angle at $w$ of the triangle $vzw$; it is opposite the side $vz$ of length $1$, hence acute, and by the law of sines it is at most $\arcsin(1/\Delta)$. Likewise $\angle(\hat w^*,u^*)\le\arcsin(1/\Delta)$. Also $\angle(u^*,u_v)=2\theta_0$, and $\angle(u_v,u_z)\le\beta$ because the rung is good.

Let $z_1$ be a T-neighbour of $z$ and $e:=z_1-z$, a unit vector with $e\cdot u_z\in(0,\beta)$, so $\angle(e,u_z)\ge\pi/2-\arcsin\beta$. Put $\eps:=\beta+\arcsin\beta+\arcsin(1/\Delta)$. By the triangle inequality for angles, $\angle(e,\hat w)\ge\pi/2-\eps$ and $\angle(e,\hat w^*)\ge\pi/2-(2\theta_0+\eps)$, so
\[
\eta:=e\cdot\hat w\le\sin\eps,\qquad \eta^*:=e\cdot\hat w^*\le\sin(2\theta_0+\eps).
\]
Both are less than $t$: since $t\ge\tau=\cos\theta_0=\sin(\pi/2-\theta_0)$ and $0\le\eps\le2\theta_0+\eps<\pi/2-\theta_0$, it suffices that $3\theta_0+\eps<\pi/2$, and indeed $\theta_0<26.4^\circ$ and $\eps<1.2^\circ$ give $3\theta_0+\eps<80.4^\circ$. By \cref{lem:B} (with $t\le1$), $|z_1w|=|z_1w^*|=\Delta$. Now $w\neq w^*$, and $z,z_1$ are two distinct points of $\Cc(w,\Delta)\cap\Cc(w^*,\Delta)$. These two points are $2\sqrt{\Delta^2-|ww^*|^2/4}\ge\sqrt3\,\Delta>1$ apart, because $|ww^*|\le\Delta$. This contradicts $|zz_1|=1$.
\end{proof}

\subsubsection*{The count in Region II}
By (T1) every point of $S$ has T-degree at most $2$. By \cref{lem:C} the D-ends of rungs with a type-(ii) K-end have T-degree $0$, and they are distinct by (R1). Let $d$ be the number of these rungs. Summing T-degrees, $2e_T\le2(s-d)$. Every other rung has a K-end in $Q$ or of type (i), and each K-vertex has at most one rung (R1), so $\rho\le d+p+5$. Hence
\[
e(S)\le e_T+\rho+C_{\mathrm{bad}}\le (s-d)+(d+p+5)+C_{\mathrm{bad}}=s+p+5+C_{\mathrm{bad}}\qquad\text{in Region II.}
\]
Mixed T-edges need no separate count here.

\subsection{Region III: \texorpdfstring{$\sqrt3/2-2\beta<\tau<\sqrt3/2+3\beta$}{sqrt3/2 - 2 beta < tau < sqrt3/2 + 3 beta}}\label{sec:regionIII}
Here $\theta_0\in(26.36^\circ,32.22^\circ)$ and $t\ge\tau>0.846$. By \cref{lem:corner}, a rung K-vertex of $\Dt$-degree at least $2$ has a normal cone of angle at least $\min\{2\theta_0,\,\pi/2-\theta_0\}>2\pi/7$, because $\tau<\sqrt3/2+3\beta<\cos(\pi/7)$ and $\tau>\sqrt3/2-2\beta>\cos(3\pi/14)$;\chk{C6} so there are at most $6$ such vertices.

Let $q_1$ and $q_2$ be the numbers of K-vertices with exactly one and exactly two good rungs. Then $\rho\le q_1+2q_2$, and $p\ge q_1+q_2-6$. It therefore suffices to prove
\begin{equation}\label{eq:eTgoal}
e_T\le s-q_2+C,
\end{equation}
since then $e(S)\le e_T+\rho+C_{\mathrm{bad}}\le s+q_1+q_2+C+C_{\mathrm{bad}}\le s+p+C+C_{\mathrm{bad}}+6$. A K-vertex with two rungs is called a \emph{two-rung vertex}; its rungs are $z=v+o_+(u_v)$ and $z'=v+o_-(u_v)$, and $z-z'=2\sigma u_v^\perp$.

\subsubsection*{(a) The case $\tau>\sqrt3/2$}
The two positions $v+o_\pm(u_v)$ are $2\sigma<1$ apart, so $q_2=0$, and~\eqref{eq:eTgoal} holds with $C=0$ by (T1).

\subsubsection*{(b) The case $\tau<\sqrt3/2$: slot killing}
Here $2\sigma>1$; more precisely $0<2\sigma-1<0.067$ and $\sigma<0.5332$.\chk{III}

\begin{lemma}\label{lem:slotkill}
Let $v$ be a two-rung vertex with rungs $z=v+o_+(u_v)$ and $z'=v+o_-(u_v)$. Then $z'$ has no T-neighbour \textup{(}pure or mixed\textup{)} on its right, and $z$ has no T-neighbour on its left.
\end{lemma}
\begin{proof}
Let $y$ be the right T-neighbour of $z'$. Then $y-z'=\cos\eps\,u_{z'}^\perp+\sin\eps\,u_{z'}$ with $\eps\in(0,\arcsin\beta)$, and $\angle(u_v,u_{z'})\le\beta$ (good rung). So the angle $\gamma'$ between $y-z'$ and $z-z'=2\sigma u_v^\perp$ is at most $\gamma:=\beta+\arcsin\beta<0.0201$. Hence
\[
|y-z|^2=1+4\sigma^2-4\sigma\cos\gamma'\le(2\sigma-1)^2+4\sigma(1-\cos\gamma)<0.005<1.\chk{III}
\]
Also $y\neq z$, since $|yz'|=1\neq2\sigma=|zz'|$. This contradicts $\delta=1$. The statement for $z$ is symmetric.
\end{proof}

\emph{Injectivity.} Suppose $v\neq\bar v$ are two-rung vertices with the same $z'=v+o_-(u_v)=\bar v+o_-(u_{\bar v})$. Put $\bar z:=\bar v+o_+(u_{\bar v})=z'+2\sigma u_{\bar v}^\perp$. Since $\angle(u_v,u_{\bar v})\le2\beta$ (both within $\beta$ of $u_{z'}$), $|\bar z-z|\le2\sigma\cdot2\beta<1$, so $\bar z=z$.\chk{III} Then $v,\bar v$ are the two points of $\Cc(z,1)\cap\Cc(z',1)$, mirror images in the line $zz'$, which has direction $u_v^\perp$. As $v-z'=\tau u_v+\sigma u_v^\perp$, we get $\bar v-z'=-\tau u_v+\sigma u_v^\perp$, and $(\bar v-z')\cdot u_{z'}\le-\tau\cos\beta+\sigma\sin\beta<0$,\chk{III} contradicting \cref{lem:identities}(a) at $z'\in D$. The same holds for the $z$-ends. So the right slots killed at the $z'$-ends are pairwise distinct, and so are the left slots killed at the $z$-ends: \cref{lem:slotkill} kills $2q_2$ distinct slots. As $e_T$ is half the number of used slots, $e_T\le\frac12(2s-2q_2)=s-q_2$, which is~\eqref{eq:eTgoal}.

\subsubsection*{(c) The case $\tau=\sqrt3/2$}
Now $\sigma=\frac12$ and $\theta_0=30^\circ$. For a two-rung vertex $v$ with $w:=w_v$, $z-z'=u_v^\perp$ has length $1$, so $vz'z$ is an equilateral triangle of side $1$, the \emph{triangle at $v$}; its \emph{base} $z'z$ is a $\delta$-edge inside $S\cap D$. By definition of rungs, $|zw|=|z'w|=\Delta$. (Property (R1) at D-vertices can fail here, since $2\sigma+2\beta>1$; the argument below does not use it.) We prove~\eqref{eq:eTgoal} with $C=C_{\mathrm{bad}}+75$.

\emph{Base determines apex.} The other point of $\Cc(z,1)\cap\Cc(z',1)$ is $\bar v=z'-\tau u_v+\frac12u_v^\perp$, and $(\bar v-z')\cdot u_{z'}\le-\tau\cos\beta+\frac12\sin\beta<0$, contradicting \cref{lem:identities}(a). So a $\delta$-edge is the base of at most one triangle, and at most $C_{\mathrm{bad}}$ triangles have a bad base. A good base joins two points of $D$, so by \cref{lem:dichotomy} it is a pure D-row T-edge; since $z-z'=u_v^\perp$ is within $\beta$ of $u_{z'}^\perp$, $z$ is the right T-neighbour of $z'$.

\emph{Small rows.} If the K-row has at most $2$ points then $q_2\le2$; if the D-row has at most $2$ points, then there is at most one base and $q_2\le1$. In both cases~\eqref{eq:eTgoal} holds with $C=2$ since $e_T\le s$. So assume both rows have at least $3$ points.

\begin{lemma}[rows are cyclic orders]\label{lem:cyclic}
If $v'$ is the right pure T-neighbour of $v$ in the K-row, then $v'$ is the successor of $v$ in the counterclockwise cyclic order of $S\setminus D$ on $\partial K$. The same holds in the D-row, on $\partial P$.
\end{lemma}
\begin{proof}
Take coordinates with $v=(0,0)$, $v'=(1,0)$. By the T-step form, $u_v=(\sin\eps,\cos\eps)$ and $u_{v'}=(-\sin\eps',\cos\eps')$ with $\eps,\eps'\in(0,\arcsin\beta)$. These are inward normals of $K$ at $v$ and $v'$, so the part of $K$ with $y\le0$ lies in the triangle $\{y\le0,\ y\ge-x\tan\eps,\ y\ge(x-1)\tan\eps'\}$. That triangle has $x\in[0,1]$ and depth at most $\tan(\arcsin\beta)$, so each of its points is within $\sqrt{1/4+\tan^2(\arcsin\beta)}<1$ of $v$ or $v'$.\chk{III} Hence the boundary arc of $K$ from $v$ to $v'$ in $\{y\le0\}$ contains no other point of the row. Traversed from $v$ to $v'$ with $K$ on the side of $u_v$ (above), this arc runs counterclockwise, because $u^\perp$ is $u$ rotated by $-90^\circ$. For the D-row, $u_z$ is an inward normal of $P$ by \cref{lem:identities}(a), and the same argument applies on $\partial P$.
\end{proof}

Let $N_K:=|S\setminus D|$ and $N_D:=|S\cap D|$, and let $B_K$ (resp.\ $B_D$) be the number of pairs $(x,\operatorname{succ}x)$ of cyclically consecutive points of the K-row (resp.\ D-row) that are not joined by a pure T-edge. By \cref{lem:cyclic} the pure T-edges of a row are among its $N$ consecutive pairs, so there are $N_K-B_K$ and $N_D-B_D$ of them. With \cref{lem:mixedT},
\begin{equation}\label{eq:eTbreaks}
e_T\le s-B_K-B_D+75.
\end{equation}
If a point $x$ has no pure right T-neighbour, then $(x,\operatorname{succ}x)$ is counted in $B_K$ (resp.\ $B_D$): otherwise $x$ would be the right T-neighbour of $\operatorname{succ}x$, and then $x=\operatorname{succ}(\operatorname{succ}x)$ by \cref{lem:cyclic}, so the row would have only $2$ points.

\emph{Rotational sense.} In both rows ``right'' means counterclockwise along the boundary (\cref{lem:cyclic}). Moreover, if $z_k\in S\cap D$ and its right T-neighbour $z_{k+1}$ both lie on $\Cc(w,\Delta)$, then $z_{k+1}$ is counterclockwise from $z_k$ about $w$. Indeed, put $\hat w_k:=(w-z_k)/\Delta$ and $e:=z_{k+1}-z_k$. Then $e\cdot\hat w_k=1/(2\Delta)$, so $e=\frac1{2\Delta}\hat w_k\pm\sqrt{1-\frac1{4\Delta^2}}\,\hat w_k^\perp$. Since $x_{z_k}$ and $w$ are both diametral partners of $z_k$ and $|x_{z_k}w|\le\Delta$, we have $\angle(u_{z_k},\hat w_k)\le\pi/3$; and $\angle(e,u_{z_k}^\perp)<\arcsin\beta$ by the T-step form. Hence $\angle(e,\hat w_k^\perp)<\pi/3+\arcsin\beta<\pi/2$, the sign is $+$, and moving along $\hat w_k^\perp$ ($\hat w_k$ rotated by $-90^\circ$) is moving counterclockwise about $w$.

\begin{lemma}[a D-row component lies on one circle]\label{lem:Bprime}
Let $Z$ be the component of the graph of pure D-row T-edges that contains a good base $z'z$ of the triangle at $v$, and let $w:=w_v$. Then $Z\subseteq\Cc(w,\Delta)$, consecutive points of $Z$ are at angle exactly $\alpha_1:=2\arcsin(1/(2\Delta))$ about $w$, all in the counterclockwise sense, and $Z$ spans an angle at most $\pi/3$ about $w$. In particular $Z$ is a path.
\end{lemma}
\begin{proof}
We know $|zw|=|z'w|=\Delta$. Let $z_k\in Z$ have both a left T-neighbour $z_{k-1}$ and a right T-neighbour $z_{k+1}$, and put $u:=u_{z_k}$. By the T-step form, $z_{k+1}-z_k=\cos\eps_1u^\perp+\sin\eps_1u$ with $\eps_1\in(0,\arcsin\beta)$, and $z_{k-1}-z_k=-\cos\eps_2u^\perp+\sin\eps_2u$ where $\sin\eps_2=(z_{k-1}-z_k)\cdot u_{z_k}$ is the value $b$ of the incoming T-edge, so $\eps_2\in(0,\arcsin\beta)$. Let $\nu=\cos\varphi\,u+\sin\varphi\,u^\perp$ be any inward normal of $P$ at $z_k$. Then $\nu\cdot(z_{k+1}-z_k)=\sin(\eps_1+\varphi)\ge0$ and $\nu\cdot(z_{k-1}-z_k)=\sin(\eps_2-\varphi)\ge0$, so $\varphi\in[-\eps_1,\eps_2]$.

Suppose $|z_kw|=\Delta$. Then $w$ is a diametral partner of $z_k$, and $\hat w_k:=(w-z_k)/\Delta$ is an inward normal of $P$ at $z_k$ (\cref{lem:identities}(a) and the remark after it). So $\eta:=(z_{k+1}-z_k)\cdot\hat w_k=\sin(\eps_1+\varphi)\in[0,\sin(2\arcsin\beta)]$, and $\sin(2\arcsin\beta)<0.0201<t$.\chk{III} By \cref{lem:B}, $|z_{k+1}w|=\Delta$. In the same way $(z_{k-1}-z_k)\cdot\hat w_k=\sin(\eps_2-\varphi)\in[0,0.0201)$ and $|z_{k-1}w|=\Delta$.

Now induct outwards from the base, with no wrap-around: going right from $z$, each vertex reached has a left T-neighbour, so whenever it has a right T-neighbour it is an interior vertex and the step applies; going left from $z'$ is symmetric. So every vertex of $Z$ lies on $\Cc(w,\Delta)$. Two points of $\Cc(w,\Delta)$ at distance $1$ are at angle $\alpha_1$ about $w$, and by the rotational-sense remark each step to the right is counterclockwise about $w$. Points of $X$ are at most $\Delta$ apart, and a chord $2\Delta\sin(\theta/2)\le\Delta$ has $\theta\le\pi/3$; so $Z$ spans at most $\pi/3$, and monotone steps cannot close up into a cycle.
\end{proof}

\begin{lemma}[one centre per component]\label{lem:W}
If $A$ and $B$ are triangles, at $v_A$ and $v_B$, whose good bases lie in the same component $Z$, then $w_{v_A}=w_{v_B}$.
\end{lemma}
\begin{proof}
Write $w_A:=w_{v_A}$, $w_B:=w_{v_B}$, and let $z'_Bz_B$ be the base of $B$. By \cref{lem:Bprime} for $A$, $z'_B,z_B\in\Cc(w_A,\Delta)$; by the definition of rungs, $z'_B,z_B\in\Cc(w_B,\Delta)$. Suppose $w_A\neq w_B$; then $w_A$ is the mirror image of $w_B$ in the line $z'_Bz_B$. Put $a:=(w_B-z_B)/\Delta$ and $e:=z'_B-z_B$. From $|z'_B-w_B|=\Delta$ we get $a\cdot e=1/(2\Delta)$ exactly. The mirror image of $a$ in the line $\R e$ is $a^*=2(a\cdot e)e-a$, and $w_A-z_B=\Delta a^*$. So $(w_A-z_B)\cdot a=\Delta(2(a\cdot e)^2-1)<0$. But $w_B$ is a diametral partner of $z_B$ and $w_A\in X$, so $(w_A-z_B)\cdot a\ge0$ by \cref{lem:identities}(a). Contradiction.
\end{proof}

So all triangles with good base in $Z$ share one centre $w$. The apex $v_A$ of such a triangle lies on $\Cc(w,\Dt)$ and on the perpendicular bisector of its base (both $v_A$ and $w$ are equidistant from the base endpoints), on the same side as the base; so, seen from $w$, it lies at the angle of the base midpoint. Order the triangles with good base in $Z$ along $Z$. Let $A,B$ be consecutive, $B$ to the right, with bases $(z_i,z_{i+1})$ and $(z_{i+k},z_{i+k+1})$. Then the apexes are at angle $k\alpha_1$ about $w$, with $k\alpha_1\le\pi/3-\alpha_1<\pi/3$ by \cref{lem:Bprime}, and $k\alpha_1\ge\alpha_2:=2\arcsin(1/(2\Dt))$ because $|v_Av_B|\ge1$.

\emph{The arc between consecutive apexes.} Let $\gamma_{AB}$ be the shorter arc of $\Cc(w,\Dt)$ from $v_A$ to $v_B$. For $x\in X\setminus\{w\}$, $\Dd(x,\Dt)\cap\Cc(w,\Dt)$ is the arc $\{\theta:\cos(\theta-\theta_x)\ge|wx|/(2\Dt)\}$ of angular width at most $\pi$. It contains $v_A,v_B\in K$, which are less than $\pi$ apart about $w$, so it contains $\gamma_{AB}$. Hence $\gamma_{AB}\subseteq K\subseteq\Dd(w,\Dt)$, and therefore $\gamma_{AB}\subseteq\partial K$.

\begin{lemma}[the K-path is forced along the arc]\label{lem:rigid}
Let $v_0:=v_A$ and follow pure right T-edges in the K-row: $v_0,v_1,\dots$. As long as $v_i\in\gamma_{AB}\setminus\{v_B\}$ and $v_i$ has a pure right T-neighbour, that neighbour is the point $v^*$ of $\gamma_{AB}$ at angle $\alpha_2$ counterclockwise from $v_i$ about $w$.
\end{lemma}
\begin{proof}
\emph{Normal fact.} $u_{v_i}=\hat w_i:=(w-v_i)/\Dt$. For $i=0$ this holds because $w=w_{v_A}$. For $v_i$ in the relative interior of $\gamma_{AB}$, $K$ contains the arc on both sides of $v_i$, and $K\subseteq\Dd(w,\Dt)$; so every inward normal $\nu$ at $v_i$ satisfies $(y-v_i)\cdot\nu\ge0$ for $y$ on the arc on both sides, which forces $\nu$ to be a positive multiple of $\hat w_i$. Since $u_{v_i}$ is an inward normal of $K$ at $v_i$, $u_{v_i}=\hat w_i$.

\emph{$v^*$ lies on $\gamma_{AB}$.} Otherwise $v_B$ would be at angle in $(0,\alpha_2)$ from $v_i$, and $|v_iv_B|<1$.

\emph{The step.} Use coordinates with second axis $u:=\hat w_i$ and first axis $u^\perp$ (counterclockwise tangent). Then $v_{i+1}-v_i=(\cos\psi,\sin\psi)$ with $\psi\in(0,\arcsin\beta)$ (T-step form, with $u_{v_i}=\hat w_i$). From $v_{i+1}\in K\subseteq\Dd(w,\Dt)$, $|v_{i+1}-w|^2=\Dt^2+1-2\Dt\sin\psi\le\Dt^2$, so $\psi\ge\psi^*:=\arcsin(1/(2\Dt))$, the elevation of $v^*$. Suppose $\psi>\psi^*$. Write $u_{v_{i+1}}=(\sin\varphi,\cos\varphi)$; $|\varphi|\le\beta$ because the T-edge is good. From $v_i\in K$, $(v_i-v_{i+1})\cdot u_{v_{i+1}}=-\sin(\psi+\varphi)\ge0$, so $\varphi\le-\psi<0$. From $v^*\in K$ (as $v^*\in\gamma_{AB}$),
\[
(v^*-v_{i+1})\cdot u_{v_{i+1}}=(\cos\psi^*-\cos\psi)\sin\varphi+(\sin\psi^*-\sin\psi)\cos\varphi\ge0 .
\]
But both terms are negative. So $\psi=\psi^*$ and $v_{i+1}=v^*$.
\end{proof}

So the path moves along $\gamma_{AB}$ in steps of exactly $\alpha_2$ about $w$, and it cannot pass $v_B$. It reaches $v_B$ only if $j\alpha_2=k\alpha_1$ for some integer $j\ge1$; otherwise it stops at a point $b_{AB}\in\gamma_{AB}\setminus\{v_B\}$ of the K-row that has no pure right T-neighbour. (Its right slot is unused or used by a mixed T-edge.) In particular a pure right T-neighbour of a point of $\gamma_{AB}$ never lies on another part of $\partial K$.

\emph{The path cannot land on $v_B$.} Since $\alpha_2>\alpha_1$, $j\alpha_2=k\alpha_1$ forces $k>j$, so $j(\alpha_2-\alpha_1)=(k-j)\alpha_1\ge\alpha_1$ and
\[
k\alpha_1=j\alpha_2\ge\frac{\alpha_1\alpha_2}{\alpha_2-\alpha_1}=\frac1{1/\alpha_1-1/\alpha_2}.
\]
For $\rho\ge10^4$ put $f(\rho):=1/(2\arcsin(1/(2\rho)))$. From $x\le\arcsin x\le x+x^3/(6(1-x^2))$ (the Taylor coefficients of $\arcsin$ beyond the first are at most $\frac16$) we get $\rho-1/(23\rho)\le f(\rho)\le\rho$. Hence $1/\alpha_1-1/\alpha_2=f(\Delta)-f(\Dt)\le t+1/(23\Dt)$. At $\tau=\sqrt3/2$ we have $t\le1$ and $t-\tau=(1-t^2)/(2\Dt)\le1/(2\Dt)$, so $t<0.8661$. Therefore $k\alpha_1\ge1/(t+1/(23\Dt))>1/0.8662>1.154>\pi/3$,\chk{III} contradicting $k\alpha_1<\pi/3$. So every pair $(A,B)$ of consecutive triangles in the same $Z$ yields a point $b_{AB}$ of the K-row in the half-open arc $[v_A,v_B)$ with no pure right T-neighbour, and hence a pair $(b_{AB},\operatorname{succ}b_{AB})$ counted in $B_K$.

\emph{Counting.} Distinct points $b$ give distinct pairs $(b,\operatorname{succ}b)$. The half-open arcs $[v_A,v_B)\subseteq\partial K$ are pairwise disjoint:
\begin{itemize}[leftmargin=2em]
\item \emph{Same $Z$:} they are consecutive intervals on one circle.
\item \emph{Different components $Z\neq Z'$ with the same centre $w$:} their angular ranges about $w$ are disjoint, since a point of one strictly inside the range of the other would be at angle less than $\alpha_1$, hence at distance less than $1$, from one of its points. The apex arcs lie in these ranges.
\item \emph{Different centres $w\neq w'$:} let $\gamma\subseteq\Cc(w,\Dt)$ and $\gamma'\subseteq\Cc(w',\Dt)$ be nondegenerate arcs of $\partial K$ sharing a point $x$. If $x$ were in the relative interior of $\gamma$, then near $x$ the boundary $\partial K$ coincides with $\gamma$ on both sides, and $\gamma'$ would contain a nondegenerate piece of $\partial K$ lying on both circles, forcing $w=w'$. So $x$ is an endpoint of both arcs; for the half-open arcs this means $x=v_A=v_{A'}$. The triangle at a K-vertex is unique, so $A=A'$.
\end{itemize}
Hence $B_K\ge\sum_Z(m_Z-1)$, where $m_Z$ is the number of triangles with good base in $Z$ and the sum is over components with $m_Z\ge1$. Each such $Z$ is a path (\cref{lem:Bprime}); its right end has no pure right T-neighbour and so gives a pair counted in $B_D$, and distinct $Z$ give distinct pairs. So $B_D\ge\#\{Z:m_Z\ge1\}$. Adding, $B_K+B_D\ge\sum_Zm_Z\ge q_2-C_{\mathrm{bad}}$, and~\eqref{eq:eTbreaks} gives
\[
e_T\le s-q_2+C_{\mathrm{bad}}+75,
\]
which is~\eqref{eq:eTgoal}.

\subsubsection*{Conclusion for Region III}
In all three cases, $e(S)\le s+p+2C_{\mathrm{bad}}+81$.

\begin{remark}
In Regions II and III, the argument of \cref{app:lean54} (ownership and pairing, rhombi, and the weighted bound with $2p$) is no longer needed. What replaces it is the observation that a rung beyond the first-order count is always paid for, either through $p$ on the $\Dt$ side or by a missing T-edge.
\end{remark}

\section{Points without a \texorpdfstring{$\Dt$}{Delta2}-partner}\label{sec:Rpoints}

In this section we are in the nondegenerate regime with $\Dt\ge10^4$; no assumption on $\tau$ is made. We bound $\deg y+m_y$ for $y\in R$.

\begin{theorem}\label{thm:N1}
All but fewer than $7.1\cdot10^7$ points $y\in R$ satisfy $m_y\le4$, and hence $\deg y+m_y\le10$.
\end{theorem}

\begin{theorem}\label{thm:N2}
Let $0<\eps_0\le1/50$, put $\eps:=\tan\eps_0$, and suppose $t>\sqrt3/2+\eps$. Then all but at most $2\pi\cdot468^2/\eps_0+5106+465^2$ points $y\in R$ satisfy $m_y\le3$, and $m_y\le2$ if $\deg y=6$; hence $\deg y+m_y\le8$. In Region~II this holds with $\eps_0=1/50$, and then the exceptional set has fewer than $7.1\cdot10^7$ points.
\end{theorem}
In Region~II, $t\ge\tau\ge\sqrt3/2+3\beta$ if $t\le1$, and $t>1$ otherwise; since $\tan(1/50)<0.0201<3\beta$, the hypothesis of \cref{thm:N2} holds with $\eps_0=1/50$.

\subsection{Structure of \texorpdfstring{$R$}{R}}
\begin{enumerate}[label=(\alph*),leftmargin=2em]
\item \emph{Points of $R\cap D$.} If $y\in R\cap D$, then by \cref{lem:identities}(a) (which holds for every point of $D$) all other points of $X$ lie in an open half-plane bounded by a line through $y$. The $\delta$-neighbours of $y$ are pairwise at least $60^\circ$ apart as seen from $y$, so $\deg y\le3$ and $\deg y+m_y\le6$.
\item \emph{$\partial K\subseteq\bigcup_{w\in X}\Cc(w,\Dt)$.} A point $x\in K$ with $|x-w|<\Dt$ for all (finitely many) $w\in X$ is an interior point of $K$. Hence a point of $X\setminus D$ on $\partial K$ has a $\Dt$-partner, i.e.\ lies in $S$. With \cref{lem:ballpoly}(a): $R\setminus D\subseteq\operatorname{int}K$ and $S\setminus D=(X\setminus D)\cap\partial K$. So every $\delta$-neighbour of a point of $X$ lies in $\operatorname{int}K$ (and then in $R\setminus D$), in $S\setminus D\subseteq\partial K$, or in $D$, which is disjoint from $K$.
\item \emph{Depth.} For $y\in R\setminus D$ put $h:=\dist(y,\partial K)>0$. Every neighbour $q$ of $y$ in $S$ is outside $\operatorname{int}K$: either $q\in S\setminus D\subseteq\partial K$, or $q\in D$ with $\dist(q,K)\ge t$ (\cref{lem:ballpoly}(c)). The segment $[y,q]$ leaves $K$ at a point $b\in\partial K$, so $h\le|yb|\le1$, and $h\le1-t$ if $q\in D$. \emph{Consequently, if $h>1$ then $m_y=0$ and $\deg y+m_y\le6$.}
\end{enumerate}

From now on let $y\in R\setminus D$ with $h\le1$. Let $p_0\in\partial K$ be a nearest boundary point, $n_0:=(p_0-y)/h$, which is an outer normal of $K$ at $p_0$, $H_0:=\{x:(x-p_0)\cdot n_0\le0\}\supseteq K$, and $\ell_0:=\partial H_0$.

\emph{Small $K$.} If the perimeter of $K$ is less than $464$, then $\operatorname{diam}K<232$ and, since the discs of radius $\frac12$ about points of $X\setminus D\subseteq K$ are disjoint, $|X\setminus D|\le(2\cdot232+1)^2=465^2$. We may therefore assume that the perimeter of $K$ is at least $464$, and we define
\[
\gamma_y:=\text{the closed arc of $\partial K$ of points at arc-length distance at most $232$ from $p_0$.}
\]

\subsection{Thin \texorpdfstring{$K$}{K}}
\begin{lemma}[thin $K$]\label{lem:thinK}
If some point $b\in\partial K\cap\Dd(y,1)$ does not lie on $\gamma_y$, then $K$ lies in a rectangle of length $L<400$ and width $W<9$. Hence $|X\setminus D|<5106$.
\end{lemma}
\begin{proof}
Let $\nu_b$ be an outer normal of $K$ at $b$, $H_b:=\{x:(x-b)\cdot\nu_b\le0\}\supseteq K$ and $\ell_b:=\partial H_b$. Then $\dist(y,\ell_0)=h\le1$ and $\dist(y,\ell_b)\le|yb|\le1$. Both arcs of $\partial K$ from $p_0$ to $b$ have length greater than $232$.

\emph{Case $\nu_b=n_0$.} Then $\ell_0=\ell_b$: otherwise one of $H_0,H_b$ is strictly contained in the other, and the point $p_0\in\ell_0$ or $b\in\ell_b$ lies outside the smaller half-plane, although both lie in $K$. So $[p_0,b]\subseteq K\cap\ell_0\subseteq\partial K$ is a boundary arc of length $|p_0b|\le2<232$, a contradiction.

\emph{Case $\nu_b\neq\pm n_0$.} The lines $\ell_0,\ell_b$ meet at a point $A$, and $Z:=H_0\cap H_b$ is a closed cone with apex $A$ and opening $\theta:=\pi-\angle(n_0,\nu_b)\in(0,\pi)$. The point $y\in Z$ has distance at most $1$ from both bounding lines; the set of such points of $Z$ is a rhombus whose far vertex is on the bisector at distance $1/\sin(\theta/2)$ from $A$, so $|yA|\le1/\sin(\theta/2)$. The point $p_0$ lies on the boundary ray of $Z$ in $\ell_0$ and $b$ on the other boundary ray, and $|Ap_0|\le|Ay|+|yp_0|\le1+1/\sin(\theta/2)$, and likewise for $|Ab|$. If $[p_0,b]\subseteq\partial K$ we conclude as before. Otherwise the line $p_0b$ meets $\partial K$ only in $p_0,b$; let $H^+$ be the closed half-plane bounded by it that contains $A$. Then $\partial K\cap H^+$ is one of the two arcs from $p_0$ to $b$, and $K\cap H^+$ is a convex subset of the triangle $Ap_0b$. By monotonicity of perimeter, the length of this arc is at most $|Ap_0|+|Ab|\le2+2/\sin(\theta/2)$. If $\theta\ge1^\circ$ this is less than $231.2<232$,\chk{N} a contradiction. So $\theta<1^\circ$.

\emph{Fatness step} ($\theta<1^\circ$, or $\nu_b=-n_0$, in which case $Z$ is a strip and we put $\theta:=0$). Let $e$ be the direction of the bisector of $Z$ (the direction of the strip if $\theta=0$), oriented so that the cone opens towards increasing values of the \emph{level} $x\cdot e$, and let $e^\perp$ be a unit normal; call $x\cdot e^\perp$ the \emph{height}. The cross-section of $Z$ at level $\lambda$ has length $w(\lambda)=w(\lambda_y)+2\tan(\theta/2)(\lambda-\lambda_y)$, where $\lambda_y$ is the level of $y$ and $w(\lambda_y)=(d_1+d_2)/\cos(\theta/2)\le2/\cos(\theta/2)=:w_y$, with $d_1,d_2\le1$ the distances from $y$ to the two lines. Let $[E_-,E_+]\ni\lambda_y$ be the range of levels of $K$, $L:=E_+-E_-$, which is positive since $K$ has interior points, and let $W$ be the range of heights of $K$. The cross-sections of $Z$ are nested, so $W\le w_y+2\tan(\theta/2)L$; also $L\le\operatorname{diam}K\le2\Dt$.

Pick $a,a'\in K$ at levels $E_-,E_+$, and let $\lambda_m:=(E_-+E_+)/2$. The section of $K$ at level $\lambda_m$ is a segment of positive length, with top $p_{\mathrm{top}}$ and bottom $p_{\mathrm{bot}}$. Let $n_{\mathrm{top}}=\cos\varphi\,e^\perp+\sin\varphi\,e$ be an outer normal of $K$ at $p_{\mathrm{top}}$. Writing $a-p_{\mathrm{top}}=-\frac L2e+v_ae^\perp$ and $a'-p_{\mathrm{top}}=\frac L2e+v_{a'}e^\perp$ with $|v_a|,|v_{a'}|\le W$, we have
\[
-\tfrac L2\sin\varphi+v_a\cos\varphi\le0,\qquad \tfrac L2\sin\varphi+v_{a'}\cos\varphi\le0 .
\]
If $\cos\varphi=0$, these two inequalities contradict each other, since $L>0$. If $\cos\varphi<0$, adding them gives $v_a+v_{a'}\ge0$; but the midpoint of $aa'$ lies in $K$ at level $\lambda_m$, so its height is at most that of $p_{\mathrm{top}}$, i.e.\ $v_a+v_{a'}\le0$. Hence $v_a+v_{a'}=0$, the midpoint of $aa'$ is $p_{\mathrm{top}}$, and then both inequalities are equalities, so $a,a'$ lie on the supporting line at $p_{\mathrm{top}}$; now $p_{\mathrm{bot}}=p_{\mathrm{top}}-\mu e^\perp$ with $\mu>0$ has $n_{\mathrm{top}}\cdot(p_{\mathrm{bot}}-p_{\mathrm{top}})=-\mu\cos\varphi>0$, so it violates the supporting half-plane. Therefore $\cos\varphi>0$, and the inequalities give $|\tan\varphi|\le2W/L$. The same holds at $p_{\mathrm{bot}}$ with $-e^\perp$ in place of $e^\perp$.

By (b) above, every boundary point $x$ of $K$ lies on some $\Cc(w,\Dt)$ with $w\in X$, and then $(x-w)/\Dt$ is an outer normal of $K$ at $x$ (as $K\subseteq\Dd(w,\Dt)$). Choose such $w_{\mathrm{top}}$ for $p_{\mathrm{top}}$ and $w_{\mathrm{bot}}$ for $p_{\mathrm{bot}}$, and use these normals above. Then, with $\tan\varphi^*:=2W/L$,
\[
(w_{\mathrm{bot}}-w_{\mathrm{top}})\cdot e^\perp=(p_{\mathrm{bot}}-p_{\mathrm{top}})\cdot e^\perp+\Dt\bigl(n_{\mathrm{top}}\cdot e^\perp-n_{\mathrm{bot}}\cdot e^\perp\bigr)\ge-W+2\Dt\cos\varphi^*,
\]
while $|w_{\mathrm{bot}}-w_{\mathrm{top}}|\le\Delta\le1.94\,\Dt$. Suppose $L\ge400$. Then $\tan\varphi^*\le2w_y/400+4\tan(\theta/2)$ and $W\le w_y+4\Dt\tan(\theta/2)$, so
\[
\frac{2\Dt\cos\varphi^*-W}{\Dt}\ge2\Bigl(1-\frac{\tan^2\varphi^*}{2}\Bigr)-4\tan(\theta/2)-\frac{w_y}{10^4}>1.96>1.94,\chk{N}
\]
a contradiction. Hence $L<400$ and $W\le w_y+800\tan(0.5^\circ)<9$. So $K$ lies in a rectangle of size less than $400\times9$, and the disjoint discs of radius $\frac12$ about the points of $X\setminus D$ lie in a rectangle of area less than $401\cdot10$; so $|X\setminus D|<4010/(\pi/4)<5106$.
\end{proof}
The argument uses only boundary points of $K$ and centres in $X$, so corners of $K$ and boundary points that are not in $X$ are covered.

\subsection{Big turning}
\begin{lemma}\label{lem:BT}
Let $\eps_0>0$. The number of $y\in R\setminus D$ with $h\le1$ and $\turn(\gamma_y)>\eps_0$ is less than $2\pi M/\eps_0$, where $M=468^2$. For $\eps_0=1/50$ this is less than $6.89\cdot10^7$.
\end{lemma}
\begin{proof}
If $b\in\gamma_y$ then $|b-p_0|\le232$ (a chord is at most the arc), so $|b-y|\le233$. The points $y$ are at least $1$ apart, so for fixed $b$ at most $(2\cdot233.5)^2\le468^2$ of them qualify (disjoint discs of radius $\frac12$ in $\Dd(b,233.5)$). Apply \cref{lem:multiplicity} to the arcs $\gamma_y$, indexed by $y$.\chk{N}
\end{proof}

\subsection{The flat case}
\begin{lemma}[flat case]\label{lem:flat}
Suppose $\partial K\cap\Dd(y,1)\subseteq\gamma_y$ and $\turn(\gamma_y)\le\eps_0$ with $0<\eps_0\le1/50$, and put $\eps:=\tan\eps_0$. Use coordinates with origin $p_0$, first axis a unit tangent $t_0\perp n_0$ and second axis $n_0$; so $y=(0,-h)$. Then there is a concave function $f$ on $[-231,231]$ with $f(0)=0$, $f\le0$ and $|f(x)|\le\eps|x|$, whose graph is contained in $\gamma_y\subseteq\partial K$, and $U:=\Dd(y,1)\cap\{(x,z):z<f(x)\}\subseteq\operatorname{int}K$. Write a neighbour $q$ of $y$ as $q=y+(\sin\psi,\cos\psi)$, with $\psi$ measured from $n_0$. Then:
\begin{enumerate}[label=\textup{(F\arabic*)},leftmargin=3em]
\item every neighbour $q\notin\operatorname{int}K$ (in particular every neighbour in $S$, and every neighbour in $D$) has $\cos\psi\ge h+f(\sin\psi)\ge h-\eps|\sin\psi|$; hence $\cos\psi>-\eps|\sin\psi|$, i.e.\ $|\psi|<\pi/2+\eps_0$;
\item every neighbour $q\in S\setminus D$ has $\cos\psi=h+f(\sin\psi)\in[h-\eps,h]$;
\item every neighbour $q\in D$ (in $S$ or in $R$) has $\cos\psi\ge h+t+f(\sin\psi)\ge h+t-\eps$;
\item every neighbour $q\in\operatorname{int}K$ has $\cos\psi<h$.
\end{enumerate}
\end{lemma}
\begin{proof}
The normal set $N(\gamma_y)$ is an arc of $S^1$ of length at most $\eps_0$ containing $n_0$. Parametrise $\gamma_y$ by arc length $\lambda\in[-232,232]$ with $\gamma(0)=p_0$. Every one-sided tangent is perpendicular to a normal in $N(\gamma_y)$, hence within $\eps_0$ of $\pm t_0$, and after orienting, $\gamma'(\lambda)\cdot t_0\ge\cos\eps_0>0$. So $x(\lambda):=\gamma(\lambda)\cdot t_0$ is strictly increasing with $|x(\pm232)|\ge232\cos\eps_0>231$, and $\gamma_y$ contains the graph of a function $f$ on $[-231,231]$ with $|f'|\le\eps$ almost everywhere; as $f(0)=0$, $|f(x)|\le\eps|x|$. Also $f\le0$ because $K\subseteq H_0$. At each graph point the outer normals have $n_0$-component at least $\cos\eps_0>0$, and $K$ lies in the corresponding supporting half-plane; so $K$ lies below every supporting line of the graph, and $f$ is concave. The set $U$ is convex (a disc intersected with the strict hypograph of a concave function; the $x$-range of $\Dd(y,1)$ is $[-1,1]$), it contains $y$ (as $-h<0=f(0)$), and $U\cap\partial K=\emptyset$, since $\partial K\cap\Dd(y,1)\subseteq\gamma_y$ and, $x$ being monotone along $\gamma_y$, the points of $\gamma_y\cap\Dd(y,1)$ lie on the graph, where $z=f(x)$. A connected set meeting $\operatorname{int}K$ and missing $\partial K$ lies in $\operatorname{int}K$; so $U\subseteq\operatorname{int}K$.

(F1) If $q\notin\operatorname{int}K$ then $q\notin U$, so $q_z\ge f(q_x)$, i.e.\ $\cos\psi-h\ge f(\sin\psi)$; and $h>0$. The inequality $\cos\psi+\tan\eps_0|\sin\psi|>0$ is equivalent to $\cos(|\psi|-\eps_0)>0$, i.e.\ $|\psi|<\pi/2+\eps_0$. (F2) A neighbour in $S\setminus D$ lies on $\partial K\cap\Dd(y,1)$, hence on the graph. (F3) The graph point $g=(\sin\psi,f(\sin\psi))$ lies in $K$ and, by (F1), directly below $q$; so $t\le\dist(q,K)\le|q-g|=\cos\psi-h-f(\sin\psi)$. (F4) A neighbour in $\operatorname{int}K$ is not on $\partial K$. If it had $\cos\psi\ge h$, it would satisfy $z\ge0\ge f(x)$; points strictly above a graph point are outside $K$ (which lies below the supporting lines), and points on the graph are on $\partial K$. So $\cos\psi-h<f(\sin\psi)\le0$.
\end{proof}
By (b) above, every neighbour of $y$ lies in $\operatorname{int}K$, in $S\setminus D$, or in $D$.

\subsection{Proofs of \texorpdfstring{\cref{thm:N1,thm:N2}}{Theorems N1 and N2}}
\begin{proof}[Proof of \cref{thm:N1}]
Take $\eps_0=1/50$. Points of $R\cap D$ have $\deg y+m_y\le6$, and points of $R\setminus D$ with $h>1$ have $m_y=0$. Let $y\in R\setminus D$ with $h\le1$. If the perimeter of $K$ is less than $464$, there are at most $465^2$ such points in total. If some point of $\partial K\cap\Dd(y,1)$ is not on $\gamma_y$, then \cref{lem:thinK} applies, and there are fewer than $5106$ such points in total. If $\turn(\gamma_y)>\eps_0$, then $y$ is one of fewer than $6.89\cdot10^7$ points (\cref{lem:BT}). Otherwise \cref{lem:flat} applies: by (F1) every neighbour of $y$ in $S$ lies in the open arc $|\psi|<\pi/2+\eps_0$, of length less than $4\pi/3$. Distinct neighbours are at least $60^\circ$ apart as seen from $y$, and an open arc of length less than $240^\circ$ contains at most $4$ such points. So $m_y\le4$. The exceptional points number fewer than $6.89\cdot10^7+5106+465^2<7.1\cdot10^7$.
\end{proof}

\begin{proof}[Proof of \cref{thm:N2}]
The exceptional points are as in the proof of \cref{thm:N1}; so let $y$ be in the flat case of \cref{lem:flat}, with $h\le1$. Note $t>\sqrt3/2+\eps>\frac12+2\eps$.
\begin{enumerate}[label=(\roman*),leftmargin=2em]
\item \emph{At most one neighbour in $D$.} By (F3) a neighbour in $D$ has $\cos\psi\ge t-\eps>\sqrt3/2$, i.e.\ $|\psi|<30^\circ$; two neighbours at angular distance at least $60^\circ$ cannot both lie in this open arc.
\item \emph{At most one neighbour in $S\setminus D$ with $\psi\in[0,\pi)$, and at most one with $\psi\in(-\pi,0]$.} If $0\le\psi_1<\psi_2$ with $\psi_2-\psi_1\ge60^\circ$ are both in $S\setminus D$, then $\psi_2<90^\circ+\eps_0<91.2^\circ$ by (F1), so $(\psi_1+\psi_2)/2\in[30^\circ,61.2^\circ)$ and $(\psi_2-\psi_1)/2\in[30^\circ,45.6^\circ)$. Hence $\cos\psi_1-\cos\psi_2=2\sin\frac{\psi_1+\psi_2}2\sin\frac{\psi_2-\psi_1}2\ge\frac12>\eps$, contradicting (F2), by which both cosines lie in $[h-\eps,h]$.
\end{enumerate}
Every neighbour in $S$ lies in $S\setminus D$ or in $D$, so $m_y\le3$; and if $m_y=3$, the neighbours in $S$ are one point of $D$ at some $\psi_D$ with $|\psi_D|<30^\circ$ and two points $K_\pm$ of $S\setminus D$.
\begin{enumerate}[label=(\roman*),leftmargin=2em,resume]
\item \emph{$\deg y=6$ and $m_y=3$ is impossible.} Six neighbours pairwise at least $60^\circ$ apart are exactly $60^\circ$ apart, at $\psi_D+60^\circ k$. So $K_\pm$ occupy two of the five directions $\psi_D\pm60^\circ$, $\psi_D\pm120^\circ$, $\psi_D+180^\circ$, and by (F2) their cosines differ by at most $\eps$. We go through the ten pairs.
\begin{itemize}[leftmargin=1.5em]
\item A pair containing $\psi_D+180^\circ$ (four pairs): $\cos(\psi_D+180^\circ)=-\cos\psi_D<-\sqrt3/2<-\eps$, contradicting (F1).
\item $\{\psi_D+60^\circ,\psi_D-60^\circ\}$: $|\cos(\psi_D+60^\circ)-\cos(\psi_D-60^\circ)|=\sqrt3|\sin\psi_D|\le\eps$, so $\cos(\psi_D\pm60^\circ)\ge\frac12\cos\psi_D-\frac{\sqrt3}2|\sin\psi_D|\ge\frac12(1-\eps^2/3)-\frac\eps2\ge\frac12-\eps$. By (F2), $h\ge\frac12-\eps$, and then (F3) gives $\cos\psi_D\ge h+t-\eps\ge\frac12+t-2\eps>1$, a contradiction.
\item $\{\psi_D+120^\circ,\psi_D-60^\circ\}$ (and symmetrically $\{\psi_D-120^\circ,\psi_D+60^\circ\}$): the two cosines are $x$ and $-x$ with $x=\cos(\psi_D+120^\circ)$, and both lie in $[h-\eps,h]$, so $|x|\le\eps/2$ and $h\le\eps$. As $\psi_D+120^\circ\in(90^\circ,150^\circ)$, this forces $\psi_D+120^\circ\le90^\circ+\arcsin(\eps/2)$, so $\psi_D+60^\circ\in(30^\circ,31.2^\circ]$. The neighbour at $\psi_D+60^\circ$ is not in $S$ (as $m_y=3$), and $\cos(\psi_D+60^\circ)>\frac12>h$, so by (F4) it is not in $\operatorname{int}K$; hence it lies in $D$, and by (F3) it has $|\psi|<30^\circ$, a contradiction.
\item $\{\psi_D+120^\circ,\psi_D-120^\circ\}$: both directions have $|\psi|>90^\circ$, so their cosines are negative and $h<\eps$ by (F2). The neighbours at $\psi_D\pm60^\circ$ are not in $S$, and one of them has $30^\circ<|\psi|\le60^\circ$, so $\cos\psi\ge\frac12>h$; by (F4) it is not in $\operatorname{int}K$, so it lies in $D$, and (F3) gives $|\psi|<30^\circ$, a contradiction.
\item The two pairs on one side, $\{\psi_D+60^\circ,\psi_D+120^\circ\}$ and $\{\psi_D-60^\circ,\psi_D-120^\circ\}$, are excluded by (ii).
\end{itemize}
\end{enumerate}
Hence $\deg y=6$ implies $m_y\le2$, and $\deg y\le5$ implies $\deg y+m_y\le8$. For $\eps_0=1/50$ the exceptional set has at most $6.89\cdot10^7+5106+465^2<7.1\cdot10^7$ points.
\end{proof}

\begin{remark}[regimes]
\Cref{thm:N1} uses nothing about rungs or $\tau$, so it holds in Regions I, II, III, including $\tau=\sqrt3/2$, $\tau\le0$ and $\tau>1$. \Cref{thm:N2} needs $t>\sqrt3/2+\eps$ with $\eps>0$ fixed; in Region~III with $\tau\ge\sqrt3/2+\eta$, $0<\eta\le1/50$, it holds with $\eps_0=\arctan(\eta/2)$, but then the exceptional count grows like $1/\eta$ and is not uniform as $\eta\to0$. For $\tau\le\sqrt3/2$ we only have \cref{thm:N1}.
\end{remark}

\begin{remark}[local sharpness of \cref{thm:N1} for $\tau<\sqrt3/2$]\label{rem:sharp}
Consider a single disc, $K=\Dd(w,\Dt)$ near $y$, with $|y-w|=\Dt-h$. A neighbour at angle $\psi$ from the outward radial direction satisfies $|q-w|^2=(\Dt-h)^2+1+2(\Dt-h)\cos\psi$. So $q$ lies on $\partial K$ if and only if $\cos\psi=c_K(h):=h+(h^2-1)/(2(\Dt-h))$, and $\dist(q,K)\ge t$ if and only if $\cos\psi\ge c_D(h):=h+t+((h+t)^2-1)/(2(\Dt-h))$; moreover $c_D(0)=\tau$ and $c_D$ is increasing in $h$.\chk{N} Two such D-feasible neighbours $60^\circ$ apart need $c_D(h)\le\cos30^\circ$, so this local picture requires $\tau<\sqrt3/2$. For $\tau<\sqrt3/2$ it occurs: with $w=0$ and $\Dt=10^4$, take on $\Cc(w,\Dt)$ chords $[p_-,p_+]$ of length $2$ (every other chord position), let $y$ be the chord midpoint, $e:=(p_+-p_-)/2$, $\nu:=y/|y|$, $z_\pm:=y\pm e/2+(\sqrt3/2)\nu$ and $a_\pm:=y\pm e/2-(\sqrt3/2)\nu$. In $60$-digit decimal arithmetic (script \texttt{r4\_nb\_exact.py}) one checks that $|yp_\pm|=|yz_\pm|=|ya_\pm|=1$, all pairwise distances are at least $1$, the $z_\pm$ lie on $\Cc(w,\Dt+t)$ with $t=0.86598790\ldots$ and $\tau=\sqrt3/2-5.0\cdot10^{-5}$, $\dist(z_\pm,\Dd(w,\Dt))=t$, and $a_\pm\in\operatorname{int}\Dd(w,\Dt)$. So every such $y$ has degree $6$, two neighbours on $\partial K$ and two D-feasible neighbours, i.e.\ $\deg y+m_y=10$ as soon as $z_\pm\in S$. In this one-centre configuration the $z_\pm$ have no $\Dt$-partner; whether $\Theta(n)$ such points $y$ with $z_\pm\in S$ can occur in a global configuration is open. So local arguments of the kind above cannot improve the constant $10$ when $\tau<\sqrt3/2$.
\end{remark}
\section{The degenerate regime \texorpdfstring{$\Delta>1.94\Dt$}{Delta > 1.94 Delta2}}\label{sec:degenerate}

\begin{proposition}\label{prop:degenerate}
If $\Delta>1.94\,\Dt$ and $n\ge4$, then $\mult(\Dt)\le n+1$.
\end{proposition}

Put $\kappa:=(\Dt/\Delta)^2<1/1.94^2$. Note that this section does not use $\delta$.

\emph{Step 1: no two disjoint diametral pairs.} Suppose $(x,x')$ and $(z,z')$ are disjoint diametral pairs with midpoints $m,m'$. The four cross distances satisfy $\sum|ab|^2=2\Delta^2+4|m-m'|^2\ge2\Delta^2$, and each is either at most $\Dt$ or equal to $\Delta$. If at most one equals $\Delta$, the sum is at most $\Delta^2(1+3\kappa)<2\Delta^2$,\chk{C10} a contradiction. So at least two cross distances equal $\Delta$.
\begin{itemize}[leftmargin=2em]
\item If two disjoint cross pairs equal $\Delta$, say $|xz|=|x'z'|=\Delta$, then $z$ and $x'$ are the two points of $\Cc(x,\Delta)\cap\Cc(z',\Delta)$, so $x,z,z',x'$ form a rhombus of side $\Delta$ whose diagonals satisfy $d_1^2+d_2^2=4\Delta^2$. One diagonal is at least $\sqrt2\Delta$, a contradiction.
\item If two cross pairs sharing a point equal $\Delta$, say $|xz|=|xz'|=\Delta$, then $xzz'$ is an equilateral triangle of side $\Delta$, and $x'\in\Cc(x,\Delta)$ with $|x'z|,|x'z'|\le\Delta$ lies on the $60^\circ$ arc of $\Cc(x,\Delta)$ between $z$ and $z'$, strictly inside it. One of $|x'z|,|x'z'|$ is at least $2\Delta\sin15^\circ$, and $(2\sin15^\circ)^2=2-\sqrt3>\kappa$,\chk{C10} so that distance lies in $(\Dt,\Delta)$, which is impossible.
\end{itemize}

\emph{Step 2: the diameter graph is a star.} A graph with no two disjoint edges is a star or a triangle. A triangle is impossible for $n\ge4$: a fourth point is not diametral, so it is within $\Dt$ of all three vertices, but $\Dt^2<\Delta^2/3$ is less than the squared circumradius.\chk{C10} So $D=\{y\}\cup A$ with $A$ the set of diametral partners of $y$; put $B:=X\setminus D\subseteq K$ and $W:=\{y\}\cup A$.

\emph{Step 3: where the $\Dt$-pairs are.}
\begin{itemize}[leftmargin=2em]
\item \emph{Within $B$: none.} For $a\in A$, $K\subseteq\Dd(y,\Dt)\cap\Dd(a,\Dt)$, a lens whose bounding box has squared diagonal less than $0.24\,\Dt^2$.\chk{C10} So $\operatorname{diam}K<\Dt$.
\item \emph{Within $A$: at most one.} $A\subseteq\Cc(y,\Delta)$, and its pairwise distances are at most $\Dt$ (a $\Delta$-pair inside $A$ would create a triangle in the diameter graph). A chord of $\Cc(y,\Delta)$ of length at most $\Dt<\Delta/1.94$ subtends an angle less than $30^\circ$ at $y$. So $A$ lies in an arc of angular width less than $60^\circ$, on which chord length is increasing in angular separation; a $\Dt$-pair in $A$ must realise the largest pairwise distance, hence be the two extreme points of $A$.
\item \emph{Between $y$ and $A$: none}, as those distances equal $\Delta$.
\end{itemize}
Hence $\mult(\Dt)\le1+I$, where $I:=\#\{(b,w):b\in B,\ w\in W,\ |bw|=\Dt\}$.

\emph{Step 4: $I\le|B|+|W|$.} We may assume $K$ has nonempty interior: an open chord between two points of $K$ lies in the interior of every disc $\Dd(w,\Dt)$, so otherwise $|B|\le|K\cap X|\le1$ and the bound is trivial. For $w\in X$ let $J_w:=\Cc(w,\Dt)\cap K$.
\begin{itemize}[leftmargin=2em]
\item \emph{Each $J_w$ is a single closed arc (possibly a point).} For $x\in X\setminus\{w\}$, $\Cc(w,\Dt)\cap\Dd(x,\Dt)$ is the closed arc $\{\cos(\theta-\theta_x)\ge|wx|/(2\Dt)\}$, of angular length at most $\pi$. For $p,q\in J_w$ we have $|pq|<\Dt<2\Dt$, so $p,q$ are not antipodal and the minor arc of $\Cc(w,\Dt)$ from $p$ to $q$ lies in every $\Dd(x,\Dt)$, hence in $K$. Since $\operatorname{diam}K<\Dt$, the angular spread of $J_w$ is less than $60^\circ$, so $J_w$ is a single closed arc.
\item \emph{No point is interior to two arcs.} Suppose $b$ is a relative-interior point of a nondegenerate arc $J_w$ and $b\in\Cc(w',\Dt)$ with $w'\in X\setminus\{w\}$. The function $f(\theta)=|w+\Dt e(\theta)-w'|^2-\Dt^2=|w-w'|^2+2\Dt\,e(\theta)\cdot(w-w')$ is $\le0$ on $J_w$ (as $K\subseteq\Dd(w',\Dt)$) and vanishes at $b$, so it has an interior local maximum with value $0$. But every local maximum of this shifted cosine is its global maximum $|w-w'|^2+2\Dt|w-w'|>0$, a contradiction.
\item \emph{Consequences.} Each nondegenerate $J_w$ is a subarc of $\partial K$, and distinct nondegenerate arcs have disjoint relative interiors. So a point $b\in B$ lies on at most two nondegenerate $W$-arcs (one on each side); any further incidences of $b$ come from degenerate arcs $J_w=\{b\}$.
\item \emph{Counting.} Let $I_{\mathrm{nd}}$ be the number of incidences from nondegenerate arcs. Then $I_{\mathrm{nd}}\le|B|+\#\{b:\text{two $W$-arcs meet at }b\}$. Mapping each such $b$ to the $W$-arc that starts at $b$ in the clockwise direction is injective, so $I_{\mathrm{nd}}\le|B|+\#\{w:J_w\text{ nondegenerate}\}$. Each degenerate arc contributes at most one incidence. Hence $I\le|B|+|W|$.
\end{itemize}
Therefore $\mult(\Dt)\le1+|B|+|A|+1=n+1$, proving \cref{prop:degenerate}.


\section{Assembly of the proof of Theorem 1}\label{sec:assembly}

\begin{table}[ht]
\centering\small
\renewcommand{\arraystretch}{1.2}
\begin{tabular}{@{}llll@{}}
\toprule
regime ($\delta=1$, $\beta=1/100$) & $e(S)\le$ & $\deg y+m_y\le$ & $\min\{\mult(\Dt),\mult(\delta)\}\le$\\
\midrule
degenerate: $\Delta>1.94\Dt$, $n\ge4$ & --- & --- & $n+1$ (\cref{prop:degenerate})\\
nondegenerate, $\Dt<10^4$ & --- & --- & $3n<3N_0$ (\cref{sec:reductions})\\
I: $\tau\le\sqrt3/2-2\beta$ (incl.\ $\tau\le0$) & $s+C$ & $10$ (\cref{thm:N1}) & $\frac{15}{11}n+C$\\
III: $\sqrt3/2-2\beta<\tau<\sqrt3/2+3\beta$ & $s+p+C$ & $10$ (\cref{thm:N1}) & $\frac{15}{11}n+C$\\
II: $\sqrt3/2+3\beta\le\tau\le1$ & $s+p+C$ & $8$ (\cref{thm:N2}) & $\frac43n+C$\\
II: $\tau>1$ ($\Leftrightarrow t>1$) & $s+C$ & $8$ (\cref{thm:N2}) & $\frac43n+C$\\
\bottomrule
\end{tabular}
\caption{The regimes in the proof of Theorem 1. In the last column $C<2\cdot10^8$; the $R$-bounds hold for all but fewer than $7.1\cdot10^7$ points of $R$.}
\end{table}

The $\tau$-intervals partition $\R$, the split at $\Delta=1.94\,\Dt$ is exhaustive, and the cases $n\le3$ (or $X$ with a single distance) are trivial.

\emph{Constants.} By \cref{prop:bad}, $C_{\mathrm{bad}}<1.83\cdot10^7$. By \cref{sec:regions}, in every $\tau$-regime of the nondegenerate regime with $\Dt\ge10^4$,
\[
e(S)\le s+p+C_e,\qquad C_e:=2C_{\mathrm{bad}}+81<3.7\cdot10^7
\]
(Region~I and $\tau>1$: $C_{\mathrm{bad}}$; Region~II: $C_{\mathrm{bad}}+5$; Region~III: $C_{\mathrm{bad}}+6$ for $\tau\neq\sqrt3/2$ and $2C_{\mathrm{bad}}+81$ for $\tau=\sqrt3/2$). By \cref{thm:N1}, $\deg y+m_y\le10$ for all but $C'<7.1\cdot10^7$ points $y\in R$, and in Region~II, by \cref{thm:N2}, $\deg y+m_y\le 8$ for all but $C'<7.1\cdot10^7$ points. \Cref{lem:lpnew} gives
\begin{align*}
\min\{\mult(\Dt),\mult(\delta)\}&\le\tfrac{15}{11}n+C_e+C'\quad\text{in general, and}\\
\min\{\mult(\Dt),\mult(\delta)\}&\le\tfrac43n+C_e+C'\quad\text{in Region II,}
\end{align*}
and $C_e+C'<1.1\cdot10^8<2\cdot10^8$. The degenerate regime gives $n+1\le\frac43n+1$, and the small cases give at most $3N_0$. So Theorem~1 holds with $C_0=3N_0+2\cdot10^8$. Part (b) is the Region~II row: its hypotheses are $\Delta\le1.94\,\Dt$ and $\tau\ge\sqrt3/2+3\beta$, and if $\Dt<10^4$ the small-case bound applies. \qed

\begin{remark}\label{rem:bottleneck}
\begin{enumerate}[leftmargin=2em,label=(\alph*)]
\item \emph{The bottleneck is the $R$-term.} The estimate $e(S)\le s+p+O(1)$ is sharp enough for $\frac43$ (\cref{rem:lpvalues}); the value $\frac{15}{11}$ comes only from the bound $\deg y+m_y\le10$ of \cref{thm:N1} in Regions I and III. By \cref{rem:sharp}, this bound cannot be improved by local arguments when $\tau<\sqrt3/2$, because the exact local configuration there has $\deg y+m_y=10$; whether it can occur for $\Theta(n)$ points in a global configuration is open.
\item \emph{What remains for $9/7$.} If $\deg y+m_y\le 6$ held for all but $o(n)$ points of $R$, Region~I would give $\frac97$, but Regions II and III would still give only $\frac43$ because of the $p$-term (\cref{rem:lpvalues}). So reaching $\frac97$ with this method would also require sharpening the $\Dt$ side in Regions II and III.
\item \emph{Consistency with the lower bound.} The construction of ienjoymath~\cite{ienjoymath} for $L\ge9/7$ lies in Region~I ($t/\delta\to0$). In it $e(S)=|S|$, every point of $S$ has $\Dt$-degree $2$ or $4$ (so $Q=\emptyset$ and $\mult(\Dt)=\frac32|S|$), and every point of $R$ has $m_y=0$ and $\deg y+m_y\le6$; we checked the last statement numerically for $n=60,120,250$ (\cref{app:verification}). This is consistent with \cref{thm:N1}, which is only an upper bound; in Region~I the gap between $\frac97$ and $\frac{15}{11}$ is exactly the gap between $6$ and $10$ in the $R$-term.
\end{enumerate}
\end{remark}
\section{A layer bound and a new class of sets}\label{sec:layers}

In this section $X$ is an arbitrary finite planar set with at least two distinct distances, $\Delta>\Dt$ are its two largest distances, and we scale so that $\Dt=1$. Every distance of $X$ is at most $1$ or equal to $\Delta$. $L_1$ is the set of extreme points of $X$ (vertices of $\conv X$), $L_2$ is the set of extreme points of $X\setminus L_1$, and $m_3:=n-|L_1|-|L_2|$. Let $G$ be the graph on $L_1\cup L_2$ whose edges are the $\Dt$-pairs, let $G'$ be its $2$-core (repeatedly delete vertices of degree less than $2$), and write $L_i':=L_i\cap V(G')$ and $\deg'=\deg_{G'}$. For $p\in X$ let $h(p)$ be the number of points of $X$ at distance $\Dt$ from $p$.

\begin{lemma}[location]\label{lem:location}
\begin{enumerate}[label=(\roman*)]
\item Both endpoints of every $\Delta$-pair lie in $L_1$.
\item Every $\Dt$-pair has an endpoint in $L_1$ and both endpoints in $L_1\cup L_2$.
\end{enumerate}
In particular $G$ contains every $\Dt$-pair of $X$.
\end{lemma}
\begin{proof}
(i) If $|pq|=\Delta$ then $X\subseteq\Dd(q,\Delta)$ and $p$ lies on its boundary circle, so $p$ is an extreme point of $X$. (ii) Let $|pq|=1$. If $p\notin L_1$, then $p$ has no $\Delta$-partner by (i), so $X\subseteq\Dd(p,1)$ and $q$, lying on the boundary circle, is in $L_1$. So some endpoint, say $p$, lies in $L_1$. Suppose $q\notin L_1$. Every point $x$ with $|px|>1$ has $|px|=\Delta$ and so lies in $L_1$. Hence $X\setminus L_1\subseteq\Dd(p,1)$ with $q$ on the boundary circle, so $q$ is extreme in $X\setminus L_1$, i.e.\ $q\in L_2$.
\end{proof}

\begin{proposition}[core degree of $L_2$]\label{prop:core2}
Every $q\in L_2'$ has $\deg'(q)=2$.
\end{proposition}
\begin{proof}
Put $q=0$. Since $q$ is not extreme it has no $\Delta$-partner (\cref{lem:location}(i)), so $X\subseteq\Dd(0,1)$ and the $G$-neighbours of $q$ lie on the unit circle. Since $q\in V(G')$, $\deg'(q)\ge2$. Suppose $v_1,v_2,v_3$ are three distinct $G'$-neighbours of $q$. Each $|v_iv_j|$ is either at most $1$ or equal to $\Delta$. The argument follows the three cases of Vesztergombi's inner-point analysis~\cite[pp.~192--194]{Ve87}, with one change: we delete no points, and the contradiction in Cases~1 and~2 is with membership in the $2$-core, via $\deg_G\ge\deg'$.

\emph{Case 3: at most one pair $v_iv_j$ at distance $\Delta$.} Label so that $|v_1v_2|,|v_2v_3|\le1$. Then $v_2$ lies on the shorter arc between $v_1$ and $v_3$, and in polar angle we may take $v_2$ at $0$, $v_1$ at $a$ and $v_3$ at $-b$ with $0<a,b\le60^\circ$. Since $q$ is not extreme, the closed half-plane $\{x:x\cdot v_2\le0\}$ contains a point $t\in X$, $t\ne q$. For $t\ne0$ with $t\cdot v_2\le0$ we get $|t-v_2|^2=|t|^2-2t\cdot v_2+1>1$, so $|tv_2|=\Delta$. By symmetry assume the polar angle $\varphi$ of $t$ lies in $[90^\circ,180^\circ]$. Then the angle between $t$ and $v_3$ is $\varphi+b\in(90^\circ,240^\circ]$, so $t\cdot v_3<0$ and $|tv_3|=\Delta$. Hence $t$ lies on the perpendicular bisector of $v_2v_3$, which passes through $0$, so $\varphi=180^\circ-b/2$. The angle between $t$ and $v_1$ is then $180^\circ-b/2-a\ge90^\circ$, so $|tv_1|^2\ge|t|^2+1>1$ and $|tv_1|=\Delta$. Now $v_1,v_2,v_3$ lie on both $\Cc(0,1)$ and $\Cc(t,\Delta)$, two distinct circles, which meet in at most two points. Contradiction. (This case uses only that $q$ is not extreme.)

\emph{Case 2: exactly two pairs at $\Delta$.} Say $|v_1v_2|=|v_1v_3|=\Delta$ and $|v_2v_3|\le1$. Then $v_1$ lies on the perpendicular bisector of $v_2v_3$, and we may take $v_1=(-1,0)$, $v_{2,3}=(\cos a,\pm\sin a)$ with $0<a\le30^\circ$ and $\Delta=2\cos(a/2)$. We claim $v_1$ has no $\Dt$-neighbour $r\ne q$; then $\deg_G(v_1)=1$, contradicting $v_1\in V(G')$. Such an $r=(x,y)$ satisfies $|r-v_1|=1$ (so $|r|^2=-2x$), $|r|\le1$, and $|r-v_i|\in[0,1]\cup\{\Delta\}$ for $i=2,3$; by symmetry $y\ge0$. Write $c=\cos a$, $s_a=\sin a$.
\begin{itemize}[leftmargin=2em]
\item \emph{Both $|r-v_2|,|r-v_3|\le1$.} Impossible: $|r-v_3|^2-1=-2x(1+c)+2ys_a\ge0$, with equality only if $r=q$.
\item \emph{$|r-v_2|=\Delta$} (the case $|r-v_3|=\Delta$ is symmetric). Using $|r|^2=-2x$, this condition is the linear equation $-2x(1+c)-2s_ay=1+2c$. Intersecting with the circle $|r-v_1|=1$ gives $x=\bigl(\pm s_a\sqrt{8c+7}-3(1+c)\bigr)/(4(1+c))$, so $|r|^2=\bigl(3(1+c)\mp s_a\sqrt{8c+7}\bigr)/(2(1+c))$. The root with $+s_a\sqrt{8c+7}$ in $|r|^2$ has $|r|^2\ge3/2$, outside the disc. For the other root, $|r|^2\le1\iff1+c\le s_a\sqrt{8c+7}\iff1+c\le(1-c)(8c+7)\iff c^2\le\frac34\iff a\ge30^\circ$. So only $a=30^\circ$ remains, and there $|r-v_3|^2=3+c-s_a\sqrt{8c+7}=2$, which lies strictly between $1$ and $\Delta^2=2+\sqrt3$. Forbidden.
\item \emph{Both equal $\Delta$.} Then $r$ lies on the $x$-axis, so $r\in\{(0,0),(-2,0)\}$, and $(-2,0)$ is outside the disc.
\end{itemize}

\emph{Case 1: all three pairs at $\Delta$.} Then $v_1,v_2,v_3$ form an equilateral triangle with circumcentre $q$, and $\Delta=\sqrt3$. Let $P\ne q$ be a further $\Dt$-neighbour of some $v_i$, so $|P-v_i|=1$ and $|P|\le1$. It cannot have both other distances at most $1$: $|P-v_i|=1$ gives $|P|^2=2P\cdot v_i$, the two conditions give $P\cdot(v_i-v_j)\le0$ for $j\ne i$, and adding and using $\sum_jv_j=0$ gives $3P\cdot v_i\le0$, so $P=q$. Hence $|P-v_j|=\sqrt3$ for some $j\ne i$. Solving exactly, there are six candidate points, two for each $v_i$, each with $|P|^2=(9-\sqrt{33})/6$ and at squared distance $(7-\sqrt{33})/2\approx0.63$ from the third vertex, so each is a $\Dt$-neighbour of exactly one $v_i$. Their pairwise squared distances are $(29-5\sqrt{33})/6\approx0.046$ for three disjoint pairs (each joining candidates of different vertices), and otherwise $4/3$, $(9-\sqrt{33})/2\approx1.628$ or $(17-\sqrt{33})/6\approx1.876$, all strictly inside the forbidden interval $(1,3)$. So a set of candidates that can coexist in $X$ has at most two elements, at most two of the $v_i$ have a second $\Dt$-neighbour, and some $v_i$ has $\deg_G(v_i)=1$, contradicting $v_i\in V(G')$.

With four or more neighbours apply the argument to any three. Hence $\deg'(q)\le2$, so $\deg'(q)=2$.
\end{proof}

\begin{lemma}[opposite-side lemma]\label{lem:M}
Let $|pa|=|pb|=|br|=1$ with $r\neq p$. Suppose $a$ lies strictly on one side of the line $pb$ and $r$ does not lie strictly on that side. Then $|ar|>1$, so $|ar|=\Delta$.
\end{lemma}
\begin{proof}
Put $p=0$, $b=(1,0)$, $a=(\cos A,-\sin A)$ with $0<A<\pi$, and $r=b+(\cos F,\sin F)$ with $0\le F\le\pi$; $F\ne\pi$ as $r\ne p$. Then $|r-a|^2-1=2(1+\cos F)(1-\cos A)+2\sin F\sin A$, where the first term is positive and the second nonnegative.
\end{proof}

\begin{lemma}\label{lem:F}
Let $p$ be an extreme point of $X$ with $\Dt$-neighbours $w_1,\dots,w_k$ in angular order. Then $\deg_G(w_i)=1$ for every $3\le i\le k-2$.
\end{lemma}
\begin{proof}
Since $p$ is extreme, some line through $p$ has $X\setminus\{p\}$ strictly on one side. So the directions of the $w_j$ from $p$ are distinct, lie in an open half-plane, and any two differ by less than $180^\circ$; for $j<i<l$, $w_j$ and $w_l$ lie strictly on opposite sides of the line $pw_i$. Suppose $r\ne p$ and $|rw_i|=1$. Put $p=0$, $w_i=(1,0)$. Then $w_{i-2},w_{i-1}$ lie strictly below the $x$-axis and $w_{i+1},w_{i+2}$ strictly above; by symmetry assume $r$ is not strictly below. Write $w_{i-2}=e(-\alpha_1)$, $w_{i-1}=e(-\alpha_2)$ with $0<\alpha_2<\alpha_1<180^\circ$. By \cref{lem:M}, $|rw_{i-2}|=|rw_{i-1}|=\Delta$, so $r$ lies on the perpendicular bisector of $w_{i-2}w_{i-1}$, the line through $0$ in direction $-\gamma$ with $\gamma=(\alpha_1+\alpha_2)/2\in(0^\circ,180^\circ)$. As $r\ne0$ is not below the axis, $r=\lambda e(180^\circ-\gamma)$ with $\lambda>0$. From $|r-w_i|=1$ we get $\lambda(\lambda+2\cos\gamma)=0$, so $\lambda=-2\cos\gamma$, whence $\gamma>90^\circ$ and $\alpha_1>90^\circ$. Then $|w_{i-2}w_i|=2\sin(\alpha_1/2)>\sqrt2>1$, so $|w_{i-2}w_i|=\Delta$. Let $\theta>0$ be the direction of $w_{i+1}$; then $\alpha_1+\theta<180^\circ$, and since the chord length $2\sin(x/2)$ is strictly increasing on $[0^\circ,180^\circ]$, $|w_{i-2}w_{i+1}|=2\sin((\alpha_1+\theta)/2)>2\sin(\alpha_1/2)=\Delta$, contradicting the maximality of $\Delta$.
\end{proof}

\begin{theorem}[degree at most $4$]\label{thm:deg4}
Every $p\in L_1'$ has $\deg'(p)\le4$.
\end{theorem}
\begin{proof}
Suppose $p\in L_1'$ has five $G'$-neighbours $u_1,\dots,u_5$ in angular order. Among the $G$-neighbours of $p$, $u_3$ has at least two on each side, so $u_3=w_i$ with $3\le i\le k-2$, and $\deg_G(u_3)=1$ by \cref{lem:F}. But $u_3\in V(G')$ forces $\deg_G(u_3)\ge2$.
\end{proof}

\begin{theorem}[layer bound]\label{thm:layer}
For every finite $X\subset\R^2$ with at least two distinct distances,
\[
\mult(\Dt)\ \le\ |L_1|+|L_2|+\tfrac12\sum_{p\in L_1'}\bigl(\deg'(p)-2\bigr).
\]
\end{theorem}
\begin{proof}
By \cref{lem:location}, $\mult(\Dt)=e(G)$. Each deletion made while pruning $G$ to its $2$-core removes at most one edge, so $e(G)\le e(G')+|V(G)\setminus V(G')|$. By \cref{prop:core2}, $e(G')=\frac12\sum_{v\in V(G')}\deg'(v)=|V(G')|+\frac12\sum_{p\in L_1'}(\deg'(p)-2)$. Add the two estimates and use $|V(G)|=|L_1|+|L_2|$.
\end{proof}

\begin{corollary}\label{cor:A1}
Assume \textup{(H3)}: no point of $L_1$ has four or more points of $X$ at distance $\Dt$. (This holds, for instance, if no four points of $X$ are concyclic.) Then
\[
\mult(\Dt)\le|L_1|+|L_2|+\tfrac12\#\{p\in L_1':\deg'(p)=3\}\le\tfrac32|L_1|+|L_2|.
\]
Consequently, under \textup{(H3)}, if $|L_1|\le2m_3$ then $\Delta$ and $\Dt$ are two distinct distances each occurring at least once and at most $n$ times.
\end{corollary}
\begin{proof}
Under (H3), $2\le\deg'(p)\le h(p)\le3$ for $p\in L_1'$. The condition $\frac32|L_1|+|L_2|\le n$ is equivalent to $|L_1|\le2m_3$, and $\mult(\Delta)\le n$ by~\cite{HP34}.
\end{proof}

\begin{corollary}[unconditional form]\label{cor:uncond}
For every finite $X$ with at least two distances,
$\mult(\Dt)\le|L_1|+|L_2|+\frac12\sum_{p\in L_1}(\min\{h(p),4\}-2)^+$. Hence $\mult(\Dt)\le n$ whenever $\sum_{p\in L_1}(\min\{h(p),4\}-2)^+\le2m_3$.
\end{corollary}
\begin{proof}
Combine \cref{thm:layer} with $2\le\deg'(p)\le\min\{h(p),4\}$ on $L_1'$ (\cref{thm:deg4}).
\end{proof}

\begin{remark}\label{rem:layers}
\begin{enumerate}[leftmargin=2em,label=(\alph*)]
\item \Cref{cor:A1} covers sets not covered by~\eqref{eq:cdl13}. For example, if $|L_1|=|L_2|=0.4n$ and $m_3=0.2n$, the three quantities in~\eqref{eq:cdl13} are $1.2n$, $1.33n$ and $1.2n$, while \cref{cor:A1} gives $\mult(\Dt)\le n$ (under (H3)). Of course~\eqref{eq:cdl13} needs no hypothesis like (H3).
\item Using $\deg'\le4$, \cref{thm:layer} recovers the bound $2|L_1|+|L_2|$ of~\cite[Theorem~1.3]{CDL25}. \Cref{prop:core2} and \cref{thm:deg4} are the inputs from~\cite{Ve96} (Propositions~3 and~5--8 there, as used in~\cite{CDL25}) needed for this; we prove them directly in the $2$-core, without deleting points. We have not checked the other inputs used in~\cite{CDL25} for the first two bounds in~\eqref{eq:cdl13}.
\item In \cref{cor:A1}, the bound $\frac32|L_1|+|L_2|$ is attained under (H3) alone by the following example, which we have checked numerically for $m=8,12,16$ (floating point): a regular $m$-gon with $m$ even, with a point at distance $\Dt$ from both endpoints of every other edge (CDL's ``servers''). These examples contain concyclic quadruples; for sets with no four concyclic points see \cref{q:genericity} below.
\end{enumerate}
\end{remark}

\begin{proposition}[distinct middle neighbours]\label{prop:middle}
Let $p\ne p'$ be points of $X$ with a common $\Dt$-neighbour $w$, and suppose that each of $p,p'$ has $\Dt$-neighbours strictly on both sides of its line to $w$. Then $X$ contains four concyclic points. In particular, if no four points of $X$ are concyclic, then distinct vertices of $\conv X$ with three $\Dt$-neighbours have distinct middle $\Dt$-neighbours.
\end{proposition}
\begin{proof}
The point $p'$ is not strictly on both sides of the line $pw$, so choose a $\Dt$-neighbour $a$ of $p$ strictly on the side of $pw$ on which $p'$ does not lie strictly. \Cref{lem:M} (with $b=w$, $r=p'$) gives $|ap'|=\Delta$. Symmetrically there is a $\Dt$-neighbour $b$ of $p'$ with $|bp|=\Delta$. Now $|pb|=|ap'|=\Delta$ and $|pa|=|bp'|=\Dt$, so $p,b,a,p'$ are concyclic by the observation in \cref{sec:open} (the four points are distinct: $|pa|=\Dt\ne\Delta=|pb|$, and $a\ne p'$, $b\ne p$ since those distances are $\Delta$). For the last sentence: if $p\in L_1$ has exactly three $\Dt$-neighbours, they lie in an open half-plane at $p$, and the middle one in angular order has neighbours of $p$ strictly on both sides of its line to $p$.
\end{proof}

\section{Sets in convex position}\label{sec:convex}

A set $X$ is in \emph{convex position} if every point of $X$ is a vertex of $\conv X$; in particular no three points of $X$ are collinear. Throughout this section $|X|=n$, $m:=\lfloor n/2\rfloor$, $N:=\binom n2$, $D=D(X)$ is the number of distinct distances, and $k=k(X)$ is the number of rare distances ($1\le\mult(d)\le n$). A distance is \emph{frequent} if $\mult(d)\ge n+1$. $R_M$ denotes the vertex set of a regular $M$-gon. By Hopf--Pannwitz, $\Delta$ is always rare; Clemen, Dumitrescu and Liu proved $k\ge2$ for $n\ge5$~\cite[Theorem~1.2]{CDL25}. For $n=5$ this is best possible: $R_5$ has $k=2$.

\begin{theorem}\label{thm:convex-summary}
Let $X$ be a set of $n$ points in convex position with $6\le n\le18$. Then at least three distinct distances each occur at least once and at most $n$ times. The same holds for every even $n\ge6$.
\end{theorem}

The cases have different status:
\begin{itemize}[leftmargin=2em]
\item $n=6,7,8$: a corollary of published results (\cref{prop:lit678}).
\item even $n\ge10$: \cref{thm:even}, a parity-sharp version of the proof of~\cite[Theorem~1.2]{CDL25};
\item $n=9$: \cref{thm:nine}, from the classification of Erd\H{o}s and Fishburn~\cite{EF96};
\item $n=11,13,15,17$: \cref{thm:1113,thm:1517}, new and computer-assisted. The $n=7$ case is re-proved the same way, without the literature.
\end{itemize}
In \cref{sec:almostcoc} we show that sets in convex position with all but $w$ points on one circle have at least $(n-3w+1)/4$ rare distances once $n+w\ge91$.

\subsection{Counting and the literature}

\begin{lemma}[parity-sharp counting]\label{lem:counting}
$(D-k)(n+1)+k\le N$. Consequently $k\ge D-m+1$ whenever $D\ge1$. More precisely, if $n=2m+1$ and $D=m+K$ with $1\le K\le m+1$, then $k\ge K+1$; if $n=2m$ and $D=m+K$ with $1\le K<2m$, then $k\ge K+2$.
\end{lemma}
\begin{proof}
The first inequality holds because frequent distances contribute at least $n+1$ pairs each and rare ones at least $1$. It gives $k\ge(D(n+1)-N)/n=D-\frac{n-1}2+\frac Dn>D-m$, hence $k\ge D-m+1$. For $n=2m+1$ the right side is $K+(m+K)/(2m+1)$, with $0<(m+K)/(2m+1)\le1$. For $n=2m$ it is $K+1+K/(2m)$, with $0<K/(2m)<1$.
\end{proof}

We use the following results from the literature.
\begin{itemize}[leftmargin=2em]
\item[(Alt)] Altman~\cite{Alt63,Alt72}: $D\ge m$; for odd $n$, $D=m$ only for $R_n$ (see also~\cite[Theorem~1]{Fi95}).
\item[(Fi)] Fishburn~\cite{Fi95}: for even $n\ge8$, $D=n/2$ if and only if $X$ is $R_n$ or $n$ vertices of $R_{n+1}$. For $n=6$ the hexagons with $D=3$ are $R_6$, $R_7$ minus a vertex, and one further hexagon~\cite[Theorem~2]{Fi95}, \cite{EF95b}, with multiplicity profiles $(6,6,3)$, $(5,5,5)$ and $(6,6,3)$ respectively (as recorded in~\cite[\S2]{CDL25}; we re-checked the profiles in exact arithmetic).
\item[(EF96)] Erd\H{o}s--Fishburn~\cite{EF96}: a convex nonagon has exactly $5$ distances if and only if its vertices are $9$ vertices of $R_{10}$ or of $R_{11}$.
\item[(r$_2$)] Erd\H{o}s--Fishburn~\cite[\S3]{EF95b}, quoting Fishburn's technical report~\cite{Fi92}: the largest possible second-largest multiplicity in a convex $n$-gon is $r_2^{(n)}=1,3,5,6,7,8$ for $n=3,\dots,8$.
\end{itemize}
A set on a circle has every distance class of maximum degree at most $2$, hence of size at most $n$; so for cocircular $X$ all distances are rare and $k=D$.

\begin{proposition}[from the literature]\label{prop:lit678}
For $n=6,7,8$ in convex position, $k\ge3$.
\end{proposition}
\begin{proof}
By (r$_2$), $r_2^{(n)}\le n$ for $n\le8$, so at most one distance is frequent and $k\ge D-1$. By (Alt), $D\ge3$ for $n=6,7$ and $D\ge4$ for $n=8$. So $k\ge3$ unless $n\in\{6,7\}$ and $D=3$. For $n=7$, $D=3$ forces $R_7$ (profile $(7,7,7)$, $k=3$); for $n=6$, (Fi) gives the profiles $(6,6,3)$ and $(5,5,5)$, so $k=3$.
\end{proof}
The value $r_2^{(7)}=7$, and the other values of $r_2^{(n)}$, rest on an unrefereed technical report~\cite{Fi92} quoted in a refereed paper~\cite{EF95b}; for $n=7$ we give an independent proof in \cref{thm:1113}.

\begin{theorem}[even $n$]\label{thm:even}
Let $X$ be in convex position with $n=2m\ge6$. Then $k\ge3$. If $n\ge8$ and $k\le3$, then $D=m+1$.
\end{theorem}
\begin{proof}
By (Alt), $D\ge m$. If $D\ge m+1$, \cref{lem:counting} gives $k\ge3$, and $k\ge4$ if $D\ge m+2$. If $D=m$ and $n=6$, (Fi) gives $k=3$. If $D=m$ and $n\ge8$, then $X$ is $R_n$, with profile $(n,\dots,n,n/2)$, or $n$ vertices of $R_{n+1}$, with profile $(n-1,\dots,n-1)$; in both cases $k=m\ge4$.
\end{proof}
\Cref{thm:even} is essentially the argument of~\cite[Theorem~1.2]{CDL25}, with the counting done with parity. It adds to the literature only for even $n\ge10$.

\begin{theorem}[$n=9$]\label{thm:nine}
Every set of $9$ points in convex position has $k\ge3$.
\end{theorem}
\begin{proof}
Here $m=4$. If $D=4$, then $X=R_9$ by (Alt) and $k=4$. If $D=5$, then $X$ is cocircular by (EF96) and $k=D=5$. If $D\ge6$, \cref{lem:counting} gives $k\ge3$.
\end{proof}

\subsection{\texorpdfstring{$n=7,11,13,15,17$}{n = 7, 11, 13, 15, 17}: an ordinal reduction and SAT certificates}\label{sec:sat}

Label $X=(p_0,\dots,p_{n-1})$ in convex cyclic order, let $d_1<\dots<d_D$ be the distinct distances, and colour each pair $\{i,j\}$ with $c(ij)=t$ where $|p_ip_j|=d_t$.

\begin{lemma}[quadrilateral inequality]\label{lem:Q}
For $i<j<k<l$ (in cyclic order), $|p_ip_k|+|p_jp_l|>|p_ip_j|+|p_kp_l|$ and $|p_ip_k|+|p_jp_l|>|p_jp_k|+|p_lp_i|$.
\end{lemma}
\begin{proof}
The diagonals of the convex quadrilateral $p_ip_jp_kp_l$ meet at a point $O$ that lies on neither of the sides $p_ip_j$, $p_kp_l$. So $|p_ip_k|+|p_jp_l|=(|p_iO|+|Op_j|)+(|p_kO|+|Op_l|)>|p_ip_j|+|p_kp_l|$, by the strict triangle inequality. The other inequality is the same.
\end{proof}

\emph{Ordinal form.} Let $x_1\le x_2$ be the two diagonal colours, and $y_1\le y_2$ the colours of one pair of opposite sides. Then it is \emph{not} the case that $x_1\le y_1$ and $x_2\le y_2$, since otherwise $|p_ip_k|+|p_jp_l|\le$ the sum of the two sides. We call a colouring of the pairs of an $n$-cycle's vertex set satisfying this for all $4$-tuples and both pairs of opposite sides a \emph{Q-colouring}. This ordinal use of the quadrilateral inequality goes back to Altman~\cite{Alt63}. Abstracting a convex distance problem to a finite search in this way is the method of Mori\'c and Pritchard~\cite{MP11}, whose Fact~3.1 is \cref{lem:Q}.

Some of our certificates also use a second valid constraint (\emph{Lemma T}): if $w\ne w'$ lie on the same open boundary arc between $p_i$ and $p_j$, then $(c(iw),c(jw))\ne(c(iw'),c(jw'))$, because a point is determined by its distances to $p_i,p_j$ and a side of the line $p_ip_j$.

\begin{lemma}[reduction to profiles]\label{lem:R}
If $k\le2$, then at least $D-2$ colour classes have size at least $n+1$, and $D\le D_{\max}(n)$, the largest $D$ with $(D-2)(n+1)+2\le N$. For $n=7,11,13,15,17$, $D_{\max}=4,6,7,8,9$.
\end{lemma}

Call a Q-colouring using exactly $D$ colours, all of them, with at least $F$ classes of size at least $n+1$, an \emph{$(n,D,F)$ ordinal counterexample}. Since \cref{lem:Q} (and Lemma T) are necessary for geometric realisability, the following suffices to prove $k\ge3$ for a given $n$: \emph{for every $2\le D\le D_{\max}(n)$ there is no $(n,D,D-2)$ ordinal counterexample.} (We check the smaller $D$ even where (Alt) would exclude them, so that no citation is needed; for $D\le4$ we check the stronger statement with $F=0$, i.e.\ no Q-colouring with $D$ colours at all.) No realisation step and no floating-point arithmetic enters.

\begin{theorem}\label{thm:1113}
Every set of $n$ points in convex position with $n\in\{7,11,13\}$ has at least three rare distances. For these $n$, the ordinal quadrilateral inequality (\cref{lem:Q}) alone implies the statement.
\end{theorem}
\begin{proof}[Proof (computer-assisted)]
By \cref{lem:R} it suffices that the following instances have no ordinal counterexample:
\begin{itemize}[leftmargin=2em]
\item $n=7$: $(7,2,0)$, $(7,3,1)$, $(7,4,2)$;
\item $n=11$: $(11,2,0)$, $(11,3,0)$, $(11,4,0)$, $(11,5,3)$, $(11,6,4)$;
\item $n=13$: $(13,2,0)$, $(13,3,0)$, $(13,4,0)$, $(13,5,0)$, $(13,6,4)$, $(13,7,5)$.
\end{itemize}
Each is encoded in CNF and shown unsatisfiable, with the refutation checked as described below and in \cref{app:verification}.
\end{proof}

\subsubsection*{Encoding}
Variables $x_{e,c}$ say that pair $e$ has colour $c$; there are exactly-one constraints per pair and at-least-one constraints per colour. In the \emph{direct} encoding of \cref{lem:Q}, each dominated $4$-tuple of colours is forbidden. In the \emph{order} encoding, auxiliary literals $y_{e,t}\Leftrightarrow c(e)\ge t$ are introduced and ``$c(a)\le t\le c(b)$ and $c(a')\le u\le c(b')$'' is forbidden for both diagonal/side matchings; this is equivalent, since sorted domination holds if and only if some matching dominates. Selector variables $f_c$ imply ``class $c$ has at least $n+1$ pairs'' through sequential counters, and $\sum_cf_c\ge F$.

\subsubsection*{Certification}
\begin{itemize}[leftmargin=2em]
\item \emph{DRAT/LRAT.} For every instance with $n=11$ and $n=13$ listed above, and for $(7,4,2)$, a DRUP proof produced by Glucose~4~\cite{glucose} was verified by \texttt{drat-trim}~\cite{drattrim}, which also emitted an LRAT proof that was re-verified by \texttt{lrat-check}~\cite{lrat}. The largest certificate, for $(13,7,5)$, has $446{,}686$ lemmas (a 37\,MB DRUP proof and an 82\,MB LRAT proof). For $n=13$ all instances were certified in the Q-only encoding. For $n=11$, $(11,6,4)$ was certified in the Q-only encoding ($35{,}510$ lemmas), and $(11,2..5,\cdot)$ in the encoding with Q and Lemma~T (Q-only is covered by the Lean proof below). A truncated proof is correctly rejected. On this route, the trusted components are the CNF encoder (a short Python program, cross-validated as in \cref{app:verification}) and the proof checkers.
\item \emph{$n=7$ without SAT.} A plain backtracking enumeration of all colourings, with Q checked as soon as its pairs are assigned, finds no ordinal counterexample for $D=2,3,4$. For $D=4$ it does so with Lemma~T switched off.
\item \emph{Lean~4.} The statements for $n=7$, $11$ and $13$ are formalised in Lean~4~\cite{lean4} with Mathlib~\cite{mathlib}; see below.
\end{itemize}

\subsubsection*{$n=15$ and $n=17$}
\begin{theorem}\label{thm:1517}
Every set of $n$ points in convex position with $n\in\{15,17\}$ has at least three rare distances. For these $n$, too, the ordinal quadrilateral inequality (\cref{lem:Q}) alone implies the statement.
\end{theorem}
\begin{proof}[Proof (computer-assisted)]
Here $m=7,8$ and $D_{\max}=m+1$ (\cref{lem:R}). It suffices to refute, for $n=15$, the instances $(15,D,0)$ for $2\le D\le6$ and $(15,D,D-2)$ for $D=7,8$; and for $n=17$, the instances $(17,D,0)$ for $2\le D\le7$ and $(17,D,D-2)$ for $D=8,9$. Only \cref{lem:Q} is encoded (order encoding, without Lemma~T). The instances with $F=D-2$ are split by the pair $\{r_1,r_2\}$ of colours that are \emph{not} required to be frequent: in the case $\{r_1,r_2\}$ every other colour class is required to have at least $n+1$ pairs, the colours $r_1,r_2$ are required to be used, and the implied bound $|r_1|+|r_2|\le N-(D-2)(n+1)$ is added as a redundant constraint. Every ordinal counterexample falls under at least one of the $\binom D2$ cases. All $5+21+28$ cases for $n=15$ and all $6+28+36$ cases for $n=17$ are unsatisfiable; each refutation was produced by the solver Kissat~\cite{kissat} as a DRAT proof and verified by \texttt{drat-trim}~\cite{drattrim}.
\end{proof}

\emph{Symmetry breaking.} For $D\ge m$ most case encodings contain symmetry-breaking constraints, justified as follows. Q-colourings, and the case conditions, are invariant under the dihedral group acting on $\{0,\dots,n-1\}$: reversing the cyclic order maps diagonals to diagonals and pairs of opposite sides to pairs of opposite sides. (Reversing the order of the colours is not a symmetry, and is not used.) Let $L^*$ be the largest cyclic length $\min(j-i,\,n-j+i)$ of a pair of the largest colour $D$. After a rotation we may assume that $\{0,L^*\}$ has colour $D$; this is encoded by selector variables $s_L$ with $s_L\Rightarrow c(0,L)=D$ and ``no pair of colour $D$ has cyclic length $>L$'', and $\bigvee_L s_L$. The reflection $i\mapsto L-i\pmod n$ fixes $\{0,L\}$, preserves cyclic lengths and exchanges the pairs $\{0,L+1\}$ and $\{L,n-1\}$, so we may also require $s_L\Rightarrow c(0,L+1)\le c(L,n-1)$. As a sanity check, all solutions of two satisfiable instances with lowered thresholds were enumerated and grouped into dihedral orbits ($392$ solutions in $42$ orbits for $n=7$, $D=4$; $21{,}348$ solutions in $1{,}280$ orbits for $n=9$, $D=5$); every orbit meets the symmetry-breaking constraints.

\emph{Cost and trust.} The hard cases are those whose non-frequent pair contains the diameter colour $D$: for $n=17$, $D=9$ these took $3$ to $12$ minutes each, and the largest DRAT proof has about $124$\,MB. The other cases took seconds. Three cases ($(15,8,\{1,2\})$, $(15,7,\{3,7\})$, $(17,9,\{2,3\})$) were re-solved with CaDiCaL~\cite{kissat} and re-verified with \texttt{drat-trim}. The same pipeline re-certifies the cases $D=m+1$ for $n=7,9,11,13$. On this route the trusted components are the CNF encoder (whose Q part is the order encoding validated in \cref{app:verification}), the symmetry-breaking argument above, and \texttt{drat-trim}; the solver is not trusted. The proofs are not stored (they are regenerated deterministically by the scripts), and \cref{thm:1517} is \emph{not} part of the Lean formalisation.

\subsubsection*{The Lean formalisation}
The file \texttt{Erdos/Erdos/E132Convex.lean} proves, without \texttt{sorry},
\begin{center}
\texttt{theorem erdos132\_convex7 (p : Fin 7 $\to$ EuclideanSpace $\R$ (Fin 2)) (hp : StrictConvexPolygon p) : 3 $\le$ (rareDists p).card}
\end{center}
and the analogous \texttt{erdos132\_convex11} and \texttt{erdos132\_convex13} (file \texttt{E132Convex13.lean}) for \texttt{Fin 11} and \texttt{Fin 13}. Here \texttt{StrictConvexPolygon p} means that \emph{every} index-increasing triple $(p_i,p_j,p_k)$, $i<j<k$, is positively oriented; that is, the vertices are given as a strictly convex polygon listed counter-clockwise. \texttt{rareDists p} is the set of distances realised by pairs $i<j$ with multiplicity at most $n$. The proof formalises \cref{lem:Q} (via the intersection point of the diagonals and strict convexity of the Euclidean norm), the passage to ranks and Q-colourings, \cref{lem:R}, and the correctness of a bit-vector encoding of the finite statement ``every Q-colouring with colours below $D_{\max}$ has at least three classes of size in $[1,n]$''. This covers all $D\le D_{\max}$ at once, with unused colours allowed. The finite statement is closed by Lean's \texttt{bv\_decide}, which bit-blasts the formula with a verified bit-blaster, calls the bundled SAT solver CaDiCaL, and checks the resulting LRAT certificate with a checker that is proved correct in Lean.

\emph{Trust base.} \texttt{\#print\allowbreak{} axioms} reports, for each theorem, the standard axioms \texttt{propext}, \texttt{Classical.choice}, \texttt{Quot.sound}, plus one axiom generated by \texttt{bv\_decide} (\texttt{\dots\_native.\allowbreak bv\_decide.\allowbreak ax\_1\_14}). That axiom asserts that the verified LRAT checker \texttt{verifyBVExpr} returns \texttt{true} on the concrete formula and certificate. The value \texttt{true} is obtained by running compiled code, not by kernel reduction, so the trust base is comparable to that of \texttt{native\_decide}: it includes the Lean compiler and runtime. Neither our Python encoder, nor any external SAT solver, nor \texttt{drat-trim} is trusted on this route. An attempt to check the $n=7$ axiom statement by kernel reduction ran for more than 20 minutes without finishing and was stopped.

\emph{Unordered point sets.} The file \texttt{E132ConvexFinset.lean} proves (using only the three standard axioms) that a finite set $S$ in convex position, in the sense of Mathlib's \texttt{ConvexIndependent}, can be listed as a \texttt{StrictConvexPolygon}. The proof sorts $S\setminus\{p_0\}$ by orientation about a point $p_0\in S$; a negatively oriented increasing triple would put a point of $S$ inside a triangle spanned by others. It also proves that the rare distances of $S$ (unordered pairs) coincide with \texttt{rareDists p} for any such listing. This gives \texttt{erdos132\_convex\{7,11,13\}\_finset}: for \texttt{S : Finset} with \texttt{S.card} $=n$ and \texttt{ConvexIndependent}, $3\le$ \texttt{(rareDistsS S).card}.

\emph{The layer bound.} The results of \cref{sec:layers} are formalised in \texttt{E132Layer.lean} without \texttt{sorry}, and \texttt{\#print\allowbreak{} axioms} reports only \texttt{propext}, \texttt{Classical.choice} and \texttt{Quot.sound} (no \texttt{bv\_decide} or \texttt{native\_decide}). This covers the localisation of $\Dt$-pairs to $L_1\cup L_2$, the degree-$2$ property of $L_2$-vertices in the $2$-core, the bound $\deg\le4$ on hull vertices, Theorem~A, Corollary~A1 and its unconditional form. There $L_1$ is the set of extreme points of $X$, shown equal to the vertex set of $\operatorname{conv}X$, and the $2$-core is the largest subset with minimum degree at least $2$.

\emph{The $\frac{54}{37}$ bound.} The $\frac{54}{37}$ bound for Problem~1.6 of~\cite{CDL25}, proved in \cref{app:lean54}, is formalised in full in \texttt{E132Main*.lean}, \texttt{E132ConvexCurve.lean} and \texttt{E132Ve87.lean} (about 7{,}000 lines):
\begin{center}
\texttt{theorem erdos132\_main : $\exists$ C\textsubscript{0} : $\R$, $\forall$ X : Finset Pt, 2 $\le$ X.card $\to$ (min (mult X (dist2 X)) (mult X (minDist X)) : $\R$) $\le$ 54/37 * X.card + C\textsubscript{0}}
\end{center}
with \texttt{mult X d} the number of unordered pairs of distinct points of $X$ at distance $d$, and \texttt{dist2}, \texttt{minDist} the second-largest and smallest distance. So the formal theorem \texttt{erdos132\_main} is the $\frac{54}{37}$ bound; Theorem~1(a) is formalised separately (below). The formal proof follows the argument of \cref{app:lean54}: \cref{sec:setup,sec:bad,sec:degenerate} and Region~I as in the text, the ownership and pairing argument in Region~II, and the weighted \cref{lem:lp2} in Region~III. There is no \texttt{sorry}, and \texttt{\#print\allowbreak{} axioms} reports only \texttt{propext}, \texttt{Classical.choice} and \texttt{Quot.sound}. The formalisation also proves Vesztergombi's bound $\mult(\Dt)\le\frac32n$ and the Hopf--Pannwitz bound, so no result is used as a black box. In three places it takes a different route from the text: (a) Vesztergombi's bound is proved in the $2$-core (every degree-$4$ core vertex has two degree-$2$ core neighbours, then double counting), with no point deletion; (b) the bad-edge bounds of \cref{sec:bad} are obtained by counting over directions of outward normals instead of arc length, so \cref{lem:lenscorner} (lens-corner arc length) is not needed; (c) the degenerate-regime incidence bound uses a local charging argument, and the mixed-T bound of \cref{lem:mixedT} is proved in a weaker form (at most $1200$ instead of $71$), which only changes the constant $C_0$.

\emph{Theorem~1\textup{(a)}.} The bound $\frac{15}{11}$ is formalised in the six files \texttt{E132Main15*.lean} (about 5{,}300 lines), which build on the $\frac{54}{37}$ development without changing it:
\begin{center}
\texttt{theorem erdos132\_main15 : $\exists$ C\textsubscript{0} : $\R$, $\forall$ X : Finset Pt, 2 $\le$ X.card $\to$ (min (mult X (dist2 X)) (mult X (minDist X)) : $\R$) $\le$ 15/11 * X.card + C\textsubscript{0}}
\end{center}
There is no \texttt{sorry}; \texttt{\#print\allowbreak{} axioms} reports only \texttt{propext}, \texttt{Classical.choice} and \texttt{Quot.sound}, and the compiled modules pass \texttt{leanchecker}, which re-checks every declaration with the kernel. The set-up, the reductions, \cref{sec:bad}, Region~I, the corner lemma, (R1), \cref{lem:mixedT} and \cref{sec:degenerate} are taken from the $\frac{54}{37}$ development; \cref{lem:identity,lem:lpnew}, \cref{lem:B,lem:C}, the count in Region~II, \cref{sec:regionIII} and \cref{thm:N1} are new. The formal proof deviates from the text as follows.
\begin{enumerate}[leftmargin=2em,label=(\alph*)]
\item \emph{\cref{thm:N1} without arc length.} The arc $\gamma_y$ and the concave-graph description of \cref{lem:flat} are replaced by a Euclidean-local case split at $y\in R\setminus D$, according to the outer normals of $K$ at the points of $\partial K\cap\Dd(y,1)$: \emph{flat} (all these normals within $\frac14$ of one unit vector; then every neighbour $q$ of $y$ in $S$ has $(q-y)\cdot\nu>-\frac14$ for a fixed unit $\nu$, and since an open arc of angle $2\arccos(-\frac14)<240^\circ$ contains at most $4$ points $60^\circ$ apart, $m_y\le4$); \emph{thin} (two nearly opposite normals; the fatness step of \cref{lem:thinK} puts $K$ in a bounded box, so $|X\setminus D|$ is bounded); and \emph{big turning} (otherwise every direction between two such normals is the outer normal of a support point within distance $150$ of $y$, so $\partial K\cap\Dd(y,150)$ has turning at least $\frac14$, which happens for $O(1)$ points $y$ by \cref{lem:multiplicity}). The exceptional set has at most $10^7$ points.
\item \emph{Region~III by slots and breaks.} For $\tau<\sqrt3/2$, \eqref{eq:eTgoal} is proved by a count of free \emph{slots} (positions next to a point of $S$ where a T-neighbour could sit), using \cref{lem:slotkill}; for $\tau=\sqrt3/2$, by a count of \emph{breaks} (points of $S$ without a pure right T-neighbour) together with the circle-arc rigidity of \cref{lem:Bprime,lem:W,lem:rigid}. \Cref{lem:cyclic} is not used.
\item \emph{\Cref{lem:C} by exact reflection.} Instead of the angle estimate in the text, the formal proof decomposes the T-step exactly with respect to the reflection that exchanges the two $\Dt$-partners and applies the threshold of \cref{lem:B} to both partners; two distinct circles of radius $\Delta$ meet in at most two points.
\item In Region~II the role of $Q$ is played by the set of rung K-ends of $\Dt$-degree $1$, and the constants differ from the text; the formal constant is $C_0=5\cdot10^9$, as for the $\frac{54}{37}$ bound.
\end{enumerate}

\emph{Limitations of the formalisation.} (i) The \texttt{bv\_decide} axiom described above remains. (ii) Only the $n=7$ LRAT certificate (1.6\,MB) is stored and replayed with \texttt{bv\_check}; for $n=11,13$ the certificates (23\,MB and 97\,MB) are not stored, so compilation re-runs CaDiCaL (about 40\,s and 4\,min respectively). (iii) \Cref{lem:lenscorner} is not formalised; the formal proof of the $\frac{54}{37}$ bound does not need it (see above). (iv) \Cref{thm:N2}, and hence Theorem~1(b) (the bound $\frac43$ in Region~II), is not formalised; nor are \cref{lem:cyclic} and the arc-length form of \cref{lem:thinK,lem:BT,lem:flat}, which the formal proof of Theorem~1(a) avoids (see above).

\begin{remark}\label{rem:ordinal}
\begin{enumerate}[leftmargin=2em,label=(\alph*)]
\item The method is tight at ``three'': $(11,6,3)$ is satisfiable (e.g.\ with class sizes $(12,9,12,12,6,4)$), as is $(9,5,2)$.
\item The ordinal statement fails for $n=5,6$: $(6,3,1)$ has an ordinal solution with profile $(9,3,3)$, so at $n=6$ geometry beyond \cref{lem:Q} is needed.
\item Without the case split by the non-frequent pair, the instance $(15,8,6)$ was unresolved after more than 20 minutes of CaDiCaL; the split into $\binom82$ cases, together with the symmetry breaking, is what makes $n=15,17$ feasible. For $n=19$ the analogous computation looks feasible but expensive (hours per hard case), and it has not been done.
\item Theorem~\ref{thm:1113} for $n=7$ also follows from the literature (\cref{prop:lit678}).
\end{enumerate}
\end{remark}

\begin{remark}[related work]
Independently, a Lean project on the Lean Pool~\cite{lyfar} proves that for point sets with coordinates in certain restricted fields (e.g.\ $\mathbb Q(\sqrt2)$), $n\le2^r$ where $r$ is the number of rare distances. This gives at least three rare distances for $n\ge8$ in those coordinate fields; it neither implies nor is implied by \cref{thm:convex-summary}. Other contributions on the forum~\cite{forum132} concern the first question of Problem \#132 (at least two rare distances) for small $n$ and general sets.
\end{remark}

\subsection{Almost cocircular sets}\label{sec:almostcoc}

For sets on one circle all distances are rare. We show that a bounded number of points off the circle cannot destroy this, up to a linear loss. The two ingredients are Kneser's theorem, which recognises cocircular sets with few distances as subsets of regular polygons (Kneser's addition theorem was applied to distance sets on circles by Momihara and Shinohara~\cite{MomiharaShinohara}; \cref{lem:kneser} is a variant of this idea), and a uniform bound for the number of vertices of a regular polygon that a point off its circle can see at chord distances, which follows from a theorem of Beukers and Smyth on torsion points of plane curves.

\begin{lemma}[cocircular sets with few distances]\label{lem:kneser}
Let $A$ be a set of $q\ge3$ points on a circle $\Gamma$ with $D(A)\le q/2+K'$, where $K'\ge0$ and $q>4K'+2$. Then $A$ is contained in the vertex set of a regular $M$-gon inscribed in $\Gamma$, for some $M$ with $q\le M\le q+2K'+1$.
\end{lemma}
\begin{proof}
Identify $\Gamma$ with $\mathbb T=\R/2\pi\mathbb Z$ by the polar angle. The distance between points with angles $a,b$ determines $a-b$ up to sign, so the difference set $A-A\subseteq\mathbb T$ has $|A-A|\le2D(A)+1\le q+2K'+1$. By Kneser's theorem~\cite{Kn53} (see also~\cite{TV06}), applied in the abelian group $\mathbb T$ to the finite sets $A$ and $-A$, with $H$ the stabiliser of $A-A$ (a finite subgroup, hence cyclic of some order $M$), we have $|A-A|\ge|A+H|+|{-A}+H|-|H|=(2t-1)M$, where $A+H$ is a union of $t$ cosets of $H$. Since $q\le tM$: if $t\ge2$ then $(2t-1)M\ge\frac{2t-1}{t}q\ge\frac32q$, so $\frac32q\le q+2K'+1$, i.e.\ $q\le4K'+2$, a contradiction. Hence $t=1$, so $A$ lies in one coset of $H$, i.e.\ in the vertex set of a regular $M$-gon on $\Gamma$, and $q\le M\le|A-A|\le q+2K'+1$.
\end{proof}

Scale $\Gamma$ to the unit circle in $\mathbb C$, let $\zeta:=e^{2\pi i/M}$ and $c_j:=|1-\zeta^j|$. We say that a point $u\in\mathbb C$ \emph{hits} the vertex $\zeta^a$ if $|u-\zeta^a|=c_j$ for some $j$.

\begin{lemma}[hit lemma]\label{lem:hit}
If $u\ne0$ and $|u|\ne1$, then $u$ hits at most $23$ vertices of the regular $M$-gon $\{\zeta^a\}$, for every $M$.
\end{lemma}
\begin{proof}
Put $\lambda:=\bar u\ne0$ and $\tau:=|u|^2-1\ne0$, and consider the Laurent polynomial
\[
F(x,y):=\lambda x+\bar\lambda x^{-1}-y-y^{-1}-\tau .
\]
For $|x|=|y|=1$ we have $|u-x|^2-|1-y|^2=-F(x,y)$. So if $u$ hits $\zeta^a$ with $|u-\zeta^a|=c_j$, then $(\zeta^a,\zeta^{j})$ and $(\zeta^a,\zeta^{-j})$ are \emph{torsion points} (both coordinates roots of unity) of the curve $F=0$.

\emph{$F$ is irreducible.} Multiply by $xy$: $G:=xyF=-x\,y^2+(\lambda x^2-\tau x+\bar\lambda)\,y-x$. As a polynomial in $y$ over $\mathbb C[x]$ it is primitive ($\gcd(x,\bar\lambda)=1$), and it is not divisible by $x$ or $y$. So a nontrivial factorisation would have two factors of degree $1$ in $y$, and the discriminant
\[
(\lambda x^2-\tau x+\bar\lambda)^2-4x^2=\bigl(\lambda x^2-(\tau+2)x+\bar\lambda\bigr)\bigl(\lambda x^2-(\tau-2)x+\bar\lambda\bigr)
\]
would be a square in $\mathbb C[x]$. The two quadratic factors have no common root (their difference is $4x$, and neither vanishes at $0$), so both would be squares, i.e.\ $(\tau+2)^2=4|\lambda|^2=(\tau-2)^2$, forcing $\tau=0$. Hence $G$ is irreducible.

\emph{$F=0$ contains no torsion coset of positive dimension.} On $\{(\zeta_1t^p,\zeta_2t^q)\}$ with $(p,q)\ne(0,0)$, $F$ becomes a Laurent polynomial in $t$: if $p,q\ne0$ its constant term is $-\tau\ne0$; if $p=0$ its $t^q$-coefficient is $-\zeta_2\ne0$; if $q=0$ its $t^p$-coefficient is $\lambda\zeta_1\ne0$.

The Newton polygon of $G$ is the square with vertices $(0,1),(1,0),(2,1),(1,2)$, of area $2$. By the theorem of Beukers and Smyth~\cite{BS02}, an irreducible plane curve $f=0$ that is not a torsion coset has at most $22\,V(f)$ torsion points, where $V(f)$ is the area of the Newton polygon of $f$. Hence the curve $F=0$ has at most $44$ torsion points. Distinct hit vertices give torsion points with distinct first coordinates, and each hit gives two of them unless $\zeta^j=\zeta^{-j}$. Here $\zeta^j=1$ is impossible (it would put $u$ on $\Gamma$), and $\zeta^j=-1$ means $|u-\zeta^a|=2$, which happens for at most two vertices, since the circle of radius $2$ about $u$ is not $\Gamma$. So if $d\le2$ hits have $c_j=2$ and $h$ is the number of hits, then $2(h-d)+d\le44$, and $h\le22+d/2\le23$.
\end{proof}

A computer search (\cref{app:verification}) found that for $M\le120$ the maximum number of hits is $13$; we have not tried to optimise the constant $23$.

\begin{theorem}[almost cocircular sets]\label{thm:AC}
Let $X$ be a set of $n$ points in convex position and $\Gamma$ a circle, and put $w:=|X\setminus\Gamma|$. Suppose $n\ge3w$ and $|X\cap\Gamma|\ge4$. Then
\[
k(X)\ \ge\ \min\Bigl\{\frac{n-3w+1}{4},\ m-\frac w2-22+\frac{22}{w+1}\Bigr\}.
\]
In particular $k(X)\ge(n-3w+1)/4$ whenever $n+w\ge91$.
\end{theorem}
\begin{proof}
Put $C:=X\cap\Gamma$, $W:=X\setminus\Gamma$, $q:=|C|=n-w\ge4$ and $K':=(q-3)/4\ge0$.

\emph{Case 1: $D(C)>q/2+K'=(3q-3)/4$.} Then $D(X)\ge D(C)$, and \cref{lem:counting} gives $k\ge D(X)-m+1>\frac{3(n-w)-3}4-\frac n2+1=\frac{n-3w+1}4$.

\emph{Case 2: $D(C)\le q/2+K'$.} Since $q>4K'+2$, \cref{lem:kneser} puts $C$ inside the vertex set of a regular $M$-gon on $\Gamma$ with $M\le q+2K'+1=(3q-1)/2$.
\begin{enumerate}[leftmargin=2em,label=(\roman*)]
\item \emph{The centre $O$ of $\Gamma$ is not in $X$.} Otherwise $O$ is a vertex of $\conv X$, so $C$ lies in a closed half-plane bounded by a line through $O$, i.e.\ in a closed semicircle, which contains at most $M/2+1\le(3q+3)/4<q$ vertices of the $M$-gon. So every $u\in W$ satisfies $u\ne O$ and $u\notin\Gamma$, and \cref{lem:hit} applies to it.
\item Within $C$ every distance graph has maximum degree at most $2$, hence at most $q$ edges. So a frequent distance of $X$ that is a chord length of the $M$-gon (a \emph{chord class}) has at least $n+1-q=w+1$ pairs that meet $W$. Each such pair is either a pair $\{u,c\}$ with $u\in W$, $c\in C$, and $c$ hit by $u$, or a pair inside $W$.
\item A frequent distance that is not a chord length has all its pairs meeting $W$ (all distances within $C$ are chord lengths). For a given distance, each $u\in W$ has at most $2$ points of $\Gamma$ at that distance, so at least $n+1-2w$ of these pairs lie inside $W$.
\item By \cref{lem:hit} there are at most $23w$ pairs of the first kind in (ii), and there are at most $\binom w2$ pairs inside $W$; every pair belongs to one distance class. Writing $F_{\mathrm{ch}}$ and $F_{\mathrm{nc}}$ for the numbers of frequent chord classes and other frequent classes,
\[
(w+1)F_{\mathrm{ch}}+(n+1-2w)F_{\mathrm{nc}}\le 23w+\tbinom w2 .
\]
\item Since $n\ge3w$ we have $n+1-2w\ge w+1$, so the number $F$ of frequent distances satisfies $F\le\bigl(23w+\binom w2\bigr)/(w+1)=\frac w2+22-\frac{22}{w+1}$.
\item By (Alt), $D(X)\ge m$, so $k=D(X)-F\ge m-\frac w2-22+\frac{22}{w+1}$.
\end{enumerate}
For the last claim, $m\ge(n-1)/2$ gives $m-\frac w2-22+\frac{22}{w+1}>\frac{n-1}2-\frac w2-22\ge\frac{n-3w+1}4$ when $n+w\ge91$.
\end{proof}

\begin{remark}\label{rem:AC}
\begin{enumerate}[leftmargin=2em,label=(\alph*)]
\item For $w=o(n)$ \cref{thm:AC} gives $k\ge(\frac14-o(1))n$. So the second question of Problem \#132 has a positive answer, in a strong form, for sets in convex position that are almost cocircular. What is missing for general convex sets is a stability statement of the form ``a convex $n$-set with at most $n/2+K$ distances has all but $o(n)$ of its points on one circle''; Kneser's theorem provides this on a circle (\cref{lem:kneser}), and we know no analogue off it.
\item The loss of $3w/4$ and the factor $\frac14$ (rather than $\frac12$) come from the threshold $q>4K'+2$ in \cref{lem:kneser}.
\end{enumerate}
\end{remark}

\section{Open problems}\label{sec:open}

\subsection*{\texorpdfstring{Problem 1.6 of~\cite{CDL25}}{Problem 1.6 of CDL}}
We now know $9/7\le L\le15/11$, i.e.\ $1.2857\ldots\le L\le1.3636\ldots$. The lower bound construction~\cite{ienjoymath} is conjectured there to be optimal. In our proof the constant $\frac{15}{11}$ comes only from the $R$-term: points $y$ without a $\Dt$-partner are only known to satisfy $\deg y+m_y\le10$ (\cref{thm:N1}), and for $\tau<\sqrt3/2$ this cannot be improved by local arguments (\cref{rem:sharp}). In Region~II the bound is already $\frac43$. So an improvement of $\frac{15}{11}$ by this method needs a non-local argument for the points of $R$ in Regions I and III, for instance one showing that $\Theta(n)$ points with $\deg y+m_y=10$ cannot coexist with a large $\mult(\Dt)$. Reaching $\frac97$ would in addition require a bound on $\mult(\Dt)$ in Regions II and III that is sharper than $\frac32s-\frac12p$ (\cref{rem:bottleneck}).

\subsection*{Genericity}
One might hope that excluding four concyclic points forces $\mult(\Dt)\le n+O(1)$, which together with \cref{thm:layer} would answer the first question of Problem \#132 for such sets. This is false. For $n=4m\ge12$, start from $n/2$ points on a circle of radius $\frac12$ interleaved with $n/2$ points on a circle of radius $R$, where $R^2+R\cos(2\pi/n)=\frac34$, and perturb generically within the solution set of the distance equations. This gives convex sets with no four concyclic points and $\mult(\Dt)=\frac54n$. We have verified these examples to $50$ digits for $n=12,20,40$ (all equalities to $5\cdot10^{-51}$, with concyclicity determinants at least $10^{-8}$). A non-convex family with $\mult(\Dt)=\frac54n-\frac32$ was verified numerically for $n\le62$. In both families every distance other than $\Delta$ and $\Dt$ occurs once, so the first question of Problem \#132 still has a positive answer for them, but the second rare distance is not $\Dt$. An elementary observation, probably folklore, explains the role of concyclicity in the extremal configurations: if $|ab|=|cd|=\Delta$ and $|ac|=|bd|=\Dt$, then $a,b,c,d$ are concyclic. (The triangles $abc$ and $dcb$ are congruent. A half-turn would make $abdc$ a parallelogram with $|ad|^2+|bc|^2=2(\Delta^2+\Dt^2)>2\Delta^2$, which is impossible, so $abdc$ is an isosceles trapezoid.)
\begin{question}\label{q:genericity}
If no four points of $X$ are concyclic, is $\mult(\Dt)\le\frac54n+O(1)$?
\end{question}
\Cref{prop:middle} records one further consequence of the no-four-concyclic hypothesis; the examples above show that such local constraints do not bound $\mult(\Dt)$ by $n+O(1)$.

\subsection*{Beyond the two largest distances}
In a counterexample to the first question of Problem \#132, every distance other than $\Delta$ occurs at least $n+1$ times, and in particular $\mult(\Dt)\ge n+1$. One might hope that then the third-largest distance $\Delta_3$ is rare. This is false: there is an $8$-point set, with coordinates in $\mathbb Q(\sqrt3,\sqrt5,\sqrt{10+2\sqrt5},\sqrt{10-2\sqrt5})$, formed by a regular pentagon and three further points, whose distance multiplicities in decreasing order of distance are $2,9,9,5,2,1$; so $\mult(\Dt)=\mult(\Delta_3)=9>8$ (checked in exact arithmetic, \cref{app:verification}). In this example the fourth distance occurs $5$ times. In computer searches (rings of points, symmetric orbits, lattice and cyclotomic pools; about $3{,}500$ sets with $\mult(\Dt)>n$), some $\Delta_j$ with $j\le4$ was always rare.
\begin{question}\label{q:top3}
Let $n\ge5$ and let $X$ be an $n$-point planar set with at least four distinct distances $\Delta>\Dt>\Delta_3>\Delta_4>\dots$. Is $\min\{\mult(\Dt),\mult(\Delta_3),\mult(\Delta_4)\}\le n$?
\end{question}
A positive answer would settle the first question of Problem \#132 for all sets with at least four distinct distances.

\subsection*{Convex position}
\begin{conjecture}[Q-ORD]\label{conj:QORD}
For every $n\ge7$, every Q-colouring of the pairs of $\{0,\dots,n-1\}$ has at least three colour classes of size at most $n$.
\end{conjecture}
Q-ORD holds for $n\in\{7,11,13,15,17\}$ (\cref{thm:1113,thm:1517}), and it would imply that every set of $n\ge7$ points in convex position has at least three rare distances. (For even $n$ we have not certified it: our solver runs for $n=8,10,12$ also used Lemma~T.) Even with the case split, the hardest cases grow from about half a minute at $n=13$ to about $12$ minutes at $n=17$, so a proof for all $n$ must be human. Deletion-minimal refutations are not local ($(7,4,2)$ already needs $25$ of the $35$ quadrilaterals), so such a proof will have to be global.

For odd $n=2m+1\ge7$ (no computation needed), \cref{lem:counting} and (Alt) show that a convex set with $k\le2$ must have $D=m+1$, $k=2$ and $m-1$ frequent distances. Moreover, writing the frequent multiplicities as $2m+2+e_d$ and the two rare ones as $a,b$, the identity $a+b+\sum_de_d=m+2$ holds. (Indeed $D=m$ gives $R_n$ with $k=m\ge3$, $D\ge m+2$ gives $k\ge3$, and the identity is $(m-1)(2m+2)+\sum e_d+a+b=m(2m+1)$.) So for odd $n\ge19$ the question in convex position is about this single kind of configuration; by \cref{thm:AC} such a configuration cannot have all but $w$ of its points on one circle if $n\ge\max\{3w+11,\,91-w\}$. Beyond three, the second question of Problem \#132 asks whether $k\to\infty$; this is open even in convex position.


\appendix
\section{Computational verification}\label{app:verification}

All scripts are in the directory \texttt{scripts/} of the accompanying repository \url{https://github.com/g8r-b8/erdos132-lean}. None of the computations below is used as a proof step, except the SAT/LRAT certificates and the Lean proof for \cref{thm:1113} and the DRAT certificates for \cref{thm:1517}. The rest are independent checks of steps proved in the text. Tags \chk{Ck} in the text refer to sections of \texttt{writeup\_a2\_checks.py}; tags \chk{II}, \chk{III}, \chk{LP} and \chk{N} refer to \texttt{r4\_regionII\_checks.py}, \texttt{r4\_regionIII\_checks.py}, \texttt{r4\_lp\_checks.py}, and \texttt{r4\_nb\_exact.py} together with \texttt{r4\_nb\_referee\_checks.py}, respectively.

\subsection*{Theorem 1}
\begin{itemize}[leftmargin=2em]
\item \texttt{writeup\_a2\_checks.py}: 58 assertions in exact rationals and rational interval arithmetic, with one-sided rational bounds for $\sqrt{\cdot}$, $\pi$, $\sin$, $\cos$, $\arcsin$. Sections: C1 the LP certificates of \cref{lem:lp,lem:lp2} (now used only in \cref{app:lean54}); C2 the identity of \cref{lem:identities}(c) and the relation between $\tau$ and $t$; C3 small-angle facts for $\Dt\ge10^4$; C4 the Region~I windows; C5 the Region~II constants ((R1), the mixed T-edge count, and, for \cref{app:lean54} only, (R2), (R2-sym), the rhombus threshold and the pairing ratio); C6 the corner-lemma angles in Region~III; C7 the lens constants; C8 case (c2); C9 the bad-edge totals; C10 the degenerate-regime constants; C11 the small-case packing bound. All pass.
\item \texttt{r4\_regionII\_checks.py}: 17 assertions in exact rational arithmetic for \cref{sec:regionII}: the constants (including the type-(i) count), the LP with the $p$-term, and the identity and threshold of \cref{lem:B} (on $500$ exact rational instances and $2000$ random exact rational instances). It also runs an adversarial floating-point search against \cref{lem:C} ($2\cdot10^5$ samples), and asserts in exact arithmetic a sufficient condition for the angle estimate in the proof of \cref{lem:C} on all of Region~II (in a formulation slightly different from the one in the text, with minimum margin $0.088$ over the samples). (The Lean proof of \cref{lem:C} uses neither condition; it argues by exact reflection, see \cref{sec:sat}.)
\item \texttt{r4\_regionIII\_checks.py}: 18 assertions in exact rational arithmetic for \cref{sec:regionIII}: the constants of \cref{lem:slotkill} and of the injectivity step, the corner-lemma angles, the mixed T-edge count $<76$, the threshold used in \cref{lem:Bprime}, the bounds on $\arcsin$ and $f$ and the margin $t+1/(23\Dt)<3/\pi$ in the arithmetic step, and the thin-cap estimate of \cref{lem:cyclic}. Floating-point searches (reported, not used): B1, the slot-killing configuration, with maximum $|y-z|=0.0691$; B2, the forced step $v_{i+1}=v^*$ of \cref{lem:rigid}, with $0$ violations in $2\cdot10^5$ samples; B3, the minimum of $|j\alpha_2-k\alpha_1|\,\Dt$, which was $1.7\cdot10^{-6}>0$.
\item \texttt{r4\_lp\_checks.py}: exact dual certificates for \cref{lem:lpnew} and \cref{rem:lpvalues}, for all values of $2c$ in the remark, with and without the $p$-term; it also re-checks the certificates of \cref{lem:lp,lem:lp2}.
\item \texttt{r4\_nb\_exact.py} (exact fractions and decimals, with rigorous rational bounds for $\sin$, $\cos$, $\tan$): (1) the closed forms of $c_K,c_D$, the identity $c_D(0)=\tau$ and the monotonicity of $c_D$ used in \cref{rem:sharp}; (2) the LP values of \cref{rem:lpvalues}; (3) the constants of \cref{lem:thinK,lem:BT}: the apex arc is at most $231.19<232$, the thin-cone margin is positive, $M\le468^2$, and the count of \cref{lem:BT} for $\eps_0=1/50$ is at most $6.88\cdot10^7$; (4) the $60$-digit verification of the configuration in \cref{rem:sharp}.
\item \texttt{r4\_nb\_referee\_checks.py}: an independent set of exact checks, with rigorous Taylor bounds, of the LP steps and the constants of \cref{thm:N1,thm:N2}.
\item \texttt{r4\_nb\_search.py}: an adversarial floating-point search over local boundaries of $K$ with random corners near a point $y$ (reported, not used). No configuration violating the bounds of \cref{thm:N1,thm:N2} in their range of validity was found.
\item \texttt{r4\_nb\_ring16\_check.py}: consistency with ienjoymath's construction (floating point). For $n=60,120,250$, every point of $R$ has $m_y=0$ and $\deg y+m_y\le6$.
\item \texttt{ref2\_numeric.py}: an independent set of 40 exact assertions (second referee pass of the $\frac{54}{37}$ argument) covering the inequalities of \texttt{writeup\_a2\_checks.py}, together with an adversarial floating-point search for (R1), (R2), (R2-sym). The search is reported, not asserted. Its maxima were $0.906$, $0.897$ and $0.914$, against the proved bounds $0.908$, $0.918$ and $0.948$. (Only (R1) is used in the proof of Theorem~1; (R2) and (R2-sym) concern \cref{app:lean54}.)
\item \texttt{ref2\_ownership.py}: an exhaustive exact check of the ownership/pairing bookkeeping of \cref{app:lean54} (not used for Theorem~1), in the abstract model the argument uses: rows that are disjoint unions of paths and cycles of length $\ge3$, rungs forming a matching, (R2), (R2-sym), and no two consecutive rhombi. It covers all row structures up to isomorphism with at most $6$ vertices per row and all partial rung matchings. There were $2{,}365{,}426$ admissible configurations without rhombi, with maximum ratio exactly $5/4$, and $2{,}741{,}183$ with rhombi allowed, with maximum ratio exactly $4/3$. Mutation tests show that (R2-sym) and the rule excluding two consecutive rhombi are both needed.
\item \texttt{ref\_degen.py}, \texttt{degen\_check.py}, \texttt{ref2\_degen.py}: adversarial searches in the degenerate regime, including $5500$ configurations built so that three or more circles $\Cc(w,\Dt)$ pass through one vertex of $\partial K$. The maximum of $\mult(\Dt)-n$ found was $0$, against the proved bound $1$.
\item \texttt{ref\_checks.py}, \texttt{a2\_checks.py}: random tests of the lens lemma (about $7000$ lenses, worst ratio $0.66$) and of the exact single-centre strip in the near-resonant regime.
\end{itemize}
The $\frac{54}{37}$ argument of \cref{app:lean54}, together with the parts of \cref{sec:setup,sec:reductions,sec:bad,sec:regions,sec:degenerate} that it shares with the present proof, was checked twice more, independently, in adversarial ``referee'' passes. The new parts of the proof of Theorem~1 were checked in independent referee passes as follows: Region~II (\cref{sec:regionII}): one; Region~III with $\tau\neq\sqrt3/2$: one; Region~III with $\tau=\sqrt3/2$: two; \cref{thm:N1,thm:N2}: one. Vesztergombi's theorem~\cite{Ve87} was checked in the original: its main theorem is $\mult(\Dt)\le\frac32n$ for all $n$, with no additive constant.

\subsection*{\texorpdfstring{Section~\ref{sec:layers}}{Layer bound}}
\begin{itemize}[leftmargin=2em]
\item \texttt{va1b\_exact.py} (sympy, mpmath interval arithmetic): exact computation of the six Case~1 candidates of \cref{prop:core2} and their pairwise squared distances, with an exhaustive check that no compatible set of candidates covers all three $v_i$. It also gives a closed form and a rigorous interval check ($6000$ subintervals, $0$ failures) of the Case~2 branch $|r-v_2|=\Delta$.
\item \texttt{va1c\_sympy.py}: symbolic checks of the identity in \cref{lem:M}, of $\lambda(\lambda+2\cos\gamma)=0$, of the chord formula, and of the bisector direction in \cref{lem:F}.
\item \texttt{va1c\_local.py}: exhaustive search for the configuration of \cref{lem:F} with unit vectors from $\mathbb Z[\omega]$ (norm $\le147$, exact), $\mathbb Z[i]$ (norm $\le325$, exact) and $N$-th roots of unity ($5\le N\le60$): $0$ hits in about $3.2\cdot10^6$ configurations. A control run with weakened hypotheses does find hits.
\item \texttt{va1\_clique.py}, \texttt{va1\_cyclo.py}, \texttt{va1c\_random.py}: searches over maximal compatible sets in lattice and cyclotomic pools, plus greedy floating-point constructions. Every set was tested against \cref{lem:location}, \cref{prop:core2}, \cref{thm:deg4}, \cref{thm:layer} and \cref{cor:A1,cor:uncond}: more than $800{,}000$ sets in the first searches and about $81{,}000$ in the second, with $0$ violations.
\end{itemize}

\subsection*{\texorpdfstring{Section~\ref{sec:convex}}{Convex position}}
\begin{itemize}[leftmargin=2em]
\item \texttt{ver\_classification\_k.py}, \texttt{ver\_hexagons.py}: exact chord-index profiles for the classified extremal sets used in \cref{thm:even,thm:nine,prop:lit678}, and strict convexity and multiplicity profiles of the three $3$-distance hexagons.
\item \texttt{counting\_parity.py}: the closed forms of \cref{lem:counting} for $n<400$.
\item \texttt{b3\_sat.py}: the CNF encoder (direct and order encodings, with or without Lemma~T). \texttt{b3\_enum.py}: the standalone backtracking enumerator for $n=7$.
\item \texttt{b3\_validate.py}: encoder validation. Every $7$-subset of $R_N$, $7\le N\le16$ ($11{,}440$ sets), and every $11$-subset of $R_N$, $N\le15$, satisfies all constraints under exact chord-index ranks, as do $3000$ (for $n=7$) and $500$ (for $n=11$) random convex polygons. The $(7,3,0)$ enumeration returns exactly the $R_7$ colouring. Every satisfying assignment printed by the SAT route was re-checked against the independent constraint code of \texttt{b3\_enum.py}.
\item Four solvers (CaDiCaL, Glucose~4, MiniSat~2.2, Lingeling) agree on the $n=7$ and $n=11$ instances.
\item \texttt{b3\_drup\_check.py}: an independent forward RUP checker, used for the $n=7$, $n=9$ and $n=11$ DRUP proofs. Its self-test rejects an empty proof of a satisfiable instance, a proof with one negated lemma, and a truncated proof.
\item \texttt{b3\_drat\_trim.py}: regenerates each CNF and Glucose proof deterministically and runs \texttt{drat-trim} with LRAT output, followed by \texttt{lrat-check}. The CNF, DRUP and LRAT files are not stored in the repository because of their size.
\end{itemize}

\begin{table}[ht]
\centering
\renewcommand{\arraystretch}{1.15}
\begin{tabular}{@{}llrll@{}}
\toprule
instance $(n,D,F)$ & encoding & lemmas & \texttt{drat-trim} & \texttt{lrat-check}\\
\midrule
$(7,4,2)$ & Q & 714 & verified & verified\\
$(11,2,0)$, $(11,3,0)$, $(11,4,0)$ & Q+T & 3, 8, 56 & verified & verified\\
$(11,5,3)$ & Q+T & 524 & verified & verified\\
$(11,6,4)$ & Q & 35,510 & verified & verified\\
$(13,2,0)$, $(13,3,0)$, $(13,4,0)$, $(13,5,0)$ & Q & 2, 7, 82, 448 & verified & verified\\
$(13,6,4)$ & Q & 8,463 & verified & verified\\
$(13,7,5)$ & Q & 446,686 & verified & verified\\
\bottomrule
\end{tabular}
\caption{Certified refutations for \cref{thm:1113}.}
\end{table}

For \cref{thm:1517}:
\begin{itemize}[leftmargin=2em]
\item \texttt{r1\_sat\_split.py}: the case-split encoder (Q-only order encoding, as in \texttt{b3\_sat.py}, plus the case constraints and the symmetry-breaking constraints). \texttt{r1\_run.py}, \texttt{r1\_batch.py}, \texttt{r1\_lowD.py}: build each CNF, solve it with Kissat writing a DRAT proof, check the proof with \texttt{drat-trim}, and delete both files. A case is counted as done only when \texttt{drat-trim} reports \texttt{s VERIFIED}. All $5+21+28$ cases for $n=15$ and $6+28+36$ cases for $n=17$ were verified; solve times were at most $56$\,s for $n=15$ with symmetry breaking ($282$\,s for the case $\{7,8\}$ without it) and at most $739$\,s for $n=17$, and the DRAT proofs were at most $124$\,MB.
\item \texttt{r1\_sb\_test.py}: the orbit test of the symmetry-breaking constraints described after \cref{thm:1517}.
\end{itemize}

For \cref{sec:almostcoc}:
\begin{itemize}[leftmargin=2em]
\item \texttt{r3\_cyc\_hits.py}: the maximum number of hits (\cref{lem:hit}) for $M\le120$. Every point with at least two hits is an intersection of two circles centred at vertices with chord radii; candidates are scanned in floating point with tolerance $10^{-7}$, and the best ones are confirmed exactly in $\mathbb Q(\zeta_M)$ by arithmetic modulo the cyclotomic polynomial. The maximum found is $13$ (at $M=60,120$).
\item \texttt{r3\_almostcoc\_check.py}: floating-point sanity check of \cref{thm:AC} on regular polygons with holes plus up to $3$ adversarial points off the circle ($M\le60$); the slack against the bound was at least $24$ in every case.
\end{itemize}

\subsection*{\texorpdfstring{Section~\ref{sec:open}}{Open problems}}
\begin{itemize}[leftmargin=2em]
\item \texttt{r2\_verify.py}: $50$-digit verification of the convex examples with $\mult(\Dt)=\frac54n$ and no four concyclic points, for $n=12,20,40$. \texttt{r2\_local\_lemmas.py}: symbolic checks of identities used in the local analysis around hull vertices with three $\Dt$-neighbours.
\item \texttt{r5\_verify\_n8.py}: exact (sympy) verification of the $8$-point example with $\mult(\Dt)=\mult(\Delta_3)=9$, including its full multiplicity profile.
\end{itemize}

\emph{Lean.} The whole development (all \texttt{E132*.lean} files: the convex-position results, the layer bound, Vesztergombi's and Hopf--Pannwitz's bounds, the $\frac{54}{37}$ bound \texttt{erdos132\_main} and the $\frac{15}{11}$ bound \texttt{erdos132\_main15}) builds with \texttt{lake exe cache get} followed by \texttt{lake build}; on a laptop this takes about 7 minutes after the Mathlib cache is downloaded, of which about 4 minutes is \texttt{bv\_decide} for $n=13$ and about 1 minute is the six \texttt{E132Main15*} files. The toolchain is Lean \texttt{v4.35.0-rc2} with Mathlib.

\section{The \texorpdfstring{$54/37$}{54/37} argument formalised in Lean}\label{app:lean54}

This appendix gives the written counterpart of the formal theorem \texttt{erdos132\_main} (\cref{sec:sat}): the bound $\min\{\mult(\Dt),\mult(\delta)\}\le\frac{54}{37}n+C_0$ with $C_0=3N_0+10^8$. It uses \cref{sec:setup}, the reductions \cref{lem:VS}, \cref{lem:E} and the small cases of \cref{sec:reductions}, \cref{sec:bad}, Region~I of \cref{sec:regionI}, (R1) and \cref{lem:mixedT} and \cref{lem:corner} of \cref{sec:regions}, and \cref{sec:degenerate}. It does not use \cref{lem:identity,lem:lpnew}, \cref{lem:B,lem:C}, the Region~III arguments of \cref{sec:regionIII}, or \cref{sec:Rpoints}. In place of the $R$-term estimate it uses the trivial bound (E), $\mult(\delta)\le e(S)+6r$.

\subsection*{Linear combinations}
\begin{lemma}\label{lem:lp}
Suppose $e(S)\le k s+C$ with $0\le k<3/2$. Then $\min\{\mult(\Dt),\mult(\delta)\}\le \frac{9}{7.5-k}\,n+C$.
\end{lemma}
\begin{proof}
Put $\lambda=(6-k)/(7.5-k)\in(0,1)$. By \cref{lem:VS,lem:E},
$\min\{\mult(\Dt),\mult(\delta)\}\le\lambda\cdot\tfrac32 s+(1-\lambda)(ks+C+6r)\le \tfrac{9}{7.5-k}(s+r)+C$.
For $k=1,\frac54,\frac43$ the factor $9/(7.5-k)$ is $\frac{18}{13},\frac{36}{25},\frac{54}{37}$.\chk{C1}
\end{proof}

\begin{lemma}\label{lem:lp2}
Suppose $e(S)\le s+2p+C$. Then $\min\{\mult(\Dt),\mult(\delta)\}\le\frac75 n+C$.
\end{lemma}
\begin{proof}
By \cref{lem:VSQ,lem:E},
$\min\le\frac45\bigl(\frac32 s-\frac12p\bigr)+\frac15\bigl(s+2p+C+6r\bigr)=\frac75 s+\frac65 r+\frac C5\le\frac75n+C$.\chk{C1}
\end{proof}

Region~I gives $e(S)\le s+C_{\mathrm{bad}}$, hence $\frac{18}{13}n+C$ by \cref{lem:lp}, and so does Region~II with $\tau>1$.

\subsection*{Region II with \texorpdfstring{$\sqrt3/2+3\beta\le\tau\le1$}{sqrt3/2 + 3 beta <= tau <= 1}}
Here $t\le1$, $t\ge\tau$ and $t-\arcsin\beta\ge0.886$; there are fewer than $71$ mixed T-edges (\cref{lem:mixedT}), and each vertex has at most one rung (R1). Rungs are written $z=v+o_{\sigma'}(u_v)$ with a sign $\sigma'\in\{+,-\}$, and $|o_{\pm}(u)-o_{\pm}(u')|\le2\sigma+\angle(u,u')$ for any choice of signs. By the T-step form, for T-steps $e_1$ at $v_0$ and $e_2$ at $z_0$ in the same direction, $\angle(e_1,e_2)\le\angle(u_{v_0},u_{z_0})+\arcsin\beta$.

\begin{lemma}[(R2), proximity]\label{lem:R2}
Let $(v_0,z_0)$ be a rung, $v_1$ the right pure T-neighbour of $v_0$ and $z_1$ the right pure T-neighbour of $z_0$; put $e_1=v_1-v_0$, $e_2=z_1-z_0$. If $v_1$ has a rung $(v_1,z')$, then $z'=z_1$.
\end{lemma}
\begin{proof}
$z'-z_1=(e_1-e_2)+(o_{\sigma_1}(u_{v_1})-o_{\sigma_0}(u_{v_0}))$. Since $\angle(u_{v_0},u_{z_0})\le\beta$ (good rung) and $\angle(u_{v_0},u_{v_1})\le\beta$ (good T-edge), $|z'-z_1|\le(\beta+\arcsin\beta)+\beta+2\sigma\approx0.918<1$.\chk{C5} Distinct points of $X$ are at least $1$ apart.
\end{proof}

\begin{lemma}[(R2-sym)]\label{lem:R2sym}
In the setting of \cref{lem:R2}, if $z_1$ has a rung $(v',z_1)$, then $v'=v_1$.
\end{lemma}
\begin{proof}
$v'=z_1-o_{\sigma'}(u_{v'})$ and $v_1=z_0-o_{\sigma_0}(u_{v_0})+e_1$, so $v'-v_1=(e_2-e_1)-(o_{\sigma'}(u_{v'})-o_{\sigma_0}(u_{v_0}))$. The angle between $u_{v'}$ and $u_{v_0}$ is at most $3\beta$, as the sum of three good-edge angles along $v'\to z_1\to z_0\to v_0$. Hence $|v'-v_1|\le(\beta+\arcsin\beta)+3\beta+2\sigma\le2\sigma+6\beta\approx0.948<1$.\chk{C5}
\end{proof}

\begin{itemize}[leftmargin=2em]
\item[(R3)] \emph{Parallel rungs and rhombi.} Rungs with the same vector $c$ satisfy $c\cdot u_v=-\tau$, so $u_v$ is one of the two unit solutions, which lie at $\pm\theta_0$ from $-c$. Rungs form a matching by (R1), so their K-ends are distinct and by \cref{lem:nocommon} have distinct normals. Hence \emph{no three rungs are parallel}, and two parallel rungs have normals exactly $2\theta_0$ apart. If $v_0v_1$, $v_1z_1$, $z_1z_0$, $z_0v_0$ are all unit edges (a \emph{rhombus}), then $v_1$ and $z_0$ are the two points of $\Cc(v_0,1)\cap\Cc(z_1,1)$, so $v_0+z_1=v_1+z_0$ and the rungs $(v_0,z_0)$ and $(v_1,z_1)$ are parallel. Since the T-edge $v_0v_1$ is good, $2\theta_0=\angle(u_{v_0},u_{v_1})\le\beta$. So \emph{rhombi exist only if $\tau\ge\cos(\beta/2)$}, and $\cos(\beta/2)>\sqrt3/2+3\beta$.\chk{C5}
\end{itemize}

\emph{Ownership.} For a rung $\mathsf r=(v,z)$, walk right from $v$ along pure K-row T-edges until reaching the first vertex after $v$ that is a rung vertex, or until the walk stops. Let $k\ge1$ be the number of vertices visited before stopping, $v$ included. Define $d$ in the same way in the D-row, starting from $z$. Then $\mathsf r$ \emph{owns} these $k+d$ vertices. Every other row vertex is \emph{unowned}: it has no rung vertex to its left on its T-path or T-cycle. Ownership is well defined, since each owned vertex belongs to the nearest rung vertex at or to its left, and that rung vertex lies on exactly one rung.

\emph{Charging.} Charge to $\mathsf r$ the rung itself and the right T-edge out of each vertex it owns. Charge every other pure T-edge to its unowned left endpoint. Every rung and pure T-edge is charged exactly once, and every vertex has at most one right T-edge. Hence $\mathrm{edges}(\mathsf r)\le1+k+d$, with equality only if both walks end at a rung vertex (``both reach''), and otherwise $\mathrm{edges}(\mathsf r)\le k+d$; unowned vertices carry at most one edge each.

\emph{Both-reach gaps are $(1,1)$ or $(\ge2,\ge2)$.} If both walks reach and $k=1$, then $z_0$ has a right neighbour $z_1$ and $v_1$ is a rung vertex; by \cref{lem:R2} its rung is $(v_1,z_1)$, so $d=1$. If both reach and $d=1$, then $v_0$ has a right neighbour $v_1$ and $z_1$ is a rung vertex; by \cref{lem:R2sym} its rung is $(v_1,z_1)$, so $k=1$. A $(1,1)$ gap is a rhombus $v_0v_1z_1z_0$ (by (R1), $z_1\neq z_0$ and $v_1\neq v_0$). A cycle with a single rung vertex gives a both-reach gap with $k$ equal to the cycle length, which is at least $3$.

\emph{Ratios.} A rhombus gap has ratio edges/vertices $=3/2$; a both-reach gap with $k,d\ge2$ has ratio at most $5/4$; every other gap has ratio at most $1$. If $\tau<\cos(\beta/2)$ there are no rhombi, so $e(S)\le\frac54s+C_{\mathrm{bad}}+71$ and \cref{lem:lp} gives $\frac{36}{25}n+C$.

\emph{Pairing.} If $\tau\ge\cos(\beta/2)$, pair each rhombus rung $\mathsf r_0=(v_0,z_0)$ with its successor $\mathsf r_1=(v_1,z_1)$. The rung $\mathsf r_1$ is not itself a rhombus rung: otherwise $\mathsf r_0,\mathsf r_1,\mathsf r_2$ would be three parallel rungs (here $v_2\neq v_0$ because T-cycles have length at least $3$). The pairing is injective, because $v_0$ is the unique left T-neighbour of $v_1$. If $\mathsf r_1$ owns $m$ vertices, the pair has ratio at most
\[
\max\Bigl\{\frac{4+m}{2+m}: m\ge4\Bigr\}\cup\Bigl\{\frac{3+m}{2+m}: m\ge2\Bigr\}=\frac43 .\chk{C5}
\]
Hence $e(S)\le\frac43s+C_{\mathrm{bad}}+71$, and \cref{lem:lp} gives $\frac{54}{37}n+C$. This is the only place where $\frac{54}{37}$ is attained.

\subsection*{Region III}
By \cref{lem:corner} and the angle bound $\min\{2\theta_0,\pi/2-\theta_0\}>2\pi/7$,\chk{C6} at most $6$ K-vertices with a rung have $\Dt$-degree at least $2$; every other K-vertex with a rung lies in $Q$. A K-vertex has at most $2$ rungs, and (T1) bounds all T-edges, mixed ones included, by $s$. Hence
\[
e(S)\le s+2p+12+C_{\mathrm{bad}},
\]
and \cref{lem:lp2} gives $\frac75n+C$. (A ``hexagonal strip'' around a single centre has average $\delta$-degree $3$ locally, so this regime genuinely needs a weighted bound such as \cref{lem:lp2} for this argument.)

\subsection*{Assembly}
\begin{table}[ht]
\centering
\renewcommand{\arraystretch}{1.2}
\begin{tabular}{@{}lll@{}}
\toprule
regime ($\delta=1$, $\beta=1/100$) & bound on $e(S)$ & bound on $\min\{\mult(\Dt),\mult(\delta)\}$\\
\midrule
degenerate: $\Delta>1.94\Dt$, $n\ge4$ & --- & $n+1$ (\cref{prop:degenerate})\\
nondegenerate, $\Dt<10^4$ & --- & $3n<3N_0$ (\cref{sec:reductions})\\
I: $\tau\le\sqrt3/2-2\beta$ (incl.\ $\tau\le0$) & $s+C$ & $\frac{18}{13}n+C$\\
III: $\sqrt3/2-2\beta<\tau<\sqrt3/2+3\beta$ & $s+2p+C$ & $\frac75n+C$\\
II: $\sqrt3/2+3\beta\le\tau<\cos(\beta/2)$ & $\frac54s+C$ & $\frac{36}{25}n+C$\\
II: $\cos(\beta/2)\le\tau\le1$ & $\frac43s+C$ & $\frac{54}{37}n+C$\\
II: $\tau>1$ ($\Leftrightarrow t>1$) & $s+C$ & $\frac{18}{13}n+C$\\
\bottomrule
\end{tabular}
\caption{The regimes of the $\frac{54}{37}$ argument; here $C<10^8$.}
\end{table}
Every constant in the last column is at most $\frac{54}{37}$,\chk{C1} and $n+1\le\frac{54}{37}n+1$. With $C_{\mathrm{bad}}+71+12<10^8$ (\cref{prop:bad}),\chk{C9} the bound holds with $C_0=3N_0+10^8$.

\begin{remark}
The two least standard parts of this argument are the ownership/pairing bookkeeping and the incidence count in Step~4 of \cref{sec:degenerate}. Both were re-derived independently, and both were also checked by computer in abstract models (\cref{app:verification}). The bookkeeping check is exhaustive up to $6+6$ row vertices, and both (R2-sym) and the ``no three parallel rungs'' fact turn out to be necessary for it. For this argument the bottleneck is the near-resonant regime $\tau\ge\cos(\beta/2)$, where rhombi occur; \cref{lem:C} removes it in the proof of Theorem~1.
\end{remark}

\bibliographystyle{amsplain}
\bibliography{refs}

\end{document}
